% Appendix fragment for A93_A95_E94_E95_neutral_ABL_EP_AWF_combined.tex.
% Insert after the existing appendices and before \begin{thebibliography}:
%   % Appendix fragment for A93_A95_E94_E95_neutral_ABL_EP_AWF_combined.tex.
% Insert after the existing appendices and before \begin{thebibliography}:
%   % Appendix fragment for A93_A95_E94_E95_neutral_ABL_EP_AWF_combined.tex.
% Insert after the existing appendices and before \begin{thebibliography}:
%   % Appendix fragment for A93_A95_E94_E95_neutral_ABL_EP_AWF_combined.tex.
% Insert after the existing appendices and before \begin{thebibliography}:
%   \input{accepted_awf_appendix_A93_A95_E94_E95}
% The main note already calls \appendix; do not call it again here.
% Uses the main note's amsmath/amssymb/amsthm/bm/booktabs/natbib/cleveref
% packages, proposition environment, and \eps, \ust, \dd, \Id, \R, \Smat.
% Bibliography keys: Shih1993, Shih1995ASM, Shih1994EVM, Shih1995EVM,
% and SCI2006. All labels introduced here have the unique prefix mawf:.
% No document class, document environment, or bibliography is added.

\section{AWF with coefficients compatible with A93, A95 and E94/E95,
a budget-determined cutoff, and a nonconstant inner kinetic-energy profile}
\label{mawf:appendix}

This appendix specifies a modified analytical wall function for each of
the three outer closures A93, A95 and E94/E95
\citep{Shih1993,Shih1995ASM,Shih1994EVM,Shih1995EVM}.
The objective is to retain its exact neutral outer equilibrium while
allowing kinetic energy to vary inside the modeled wall layer.
Two auxiliary coefficients are matched to the selected outer branch;
the dissipation cutoff is determined by the reference energy budget;
and the prescribed redistribution source is removed. Ordinary diffusion
of the transported inner kinetic energy then balances the local
production--dissipation difference.

The rough-wall viscosity kernel and analytical momentum reconstruction
start from the AWF of Suga, Craft and Iacovides \citep{SCI2006}.
The compatible coefficients, endpoint sampling, budget cutoff and
variable inner reference specified below are explicit modifications.
Selecting this variant replaces the earlier wall-coefficient,
sampling and \(\chi\)-dependent scalar-source choices. Its own
dissipation law is retained with \(C_k^w=0\).

\subsection{Scope, endpoint sampling, and compatible coefficients}
\label{mawf:coefficients}

Use the selected high-Re outer equations: molecular momentum and scalar
diffusion are omitted from the outer analytical reference, and the
E94/E95 dissipation denominator is reduced from
\(k+\sqrt{\nu\eps}\) to \(k\) in that outer limit.
The inner AWF retains the fixed physical \(\nu>0\) in its momentum,
production and scalar-diffusion formulas. These are two declared
equation regimes; the result is not an exact logarithmic solution of
an unchanged finite-Re outer system.

Only inner \(k_w\) is transported. Inner dissipation is a dependent
total sink, and no inner dissipation transport equation is introduced.
The inner stress uses the auxiliary linear viscosity specified below;
the original A93/A95 nonlinear stress laws remain the outer closures.
Coordinates are \(x\) streamwise, \(y\) spanwise, and \(z\) vertical.
All stresses are per unit density.

Let \(m\in\{93,95,E\}\) identify the outer model.
Its neutral equilibrium ratio \(q_m=k_{o,0}/\ust^2\) is
\begin{equation}
 \boxed{
 q_{93}=\frac{57+\sqrt{6249}}{40},\qquad
 q_{95}=\frac{3+\sqrt{61}}{2\sqrt2},\qquad
 q_E=\frac{3+\sqrt{41}}{2\sqrt2}.}
 \label{mawf:q}
\end{equation}
These roots follow from \(C_{\mu,m}(q_m)q_m^2=1\) in simple shear.
The nonlinear corrections of A93 and A95 have zero \(xz\) component
there, but the full Reynolds shear is \(R_{xz,o,0}=-\ust^2\).

The outer dissipation equation additionally requires
\begin{equation}
 \boxed{\kappa^2q_m=\sigma_{\eps,m}\Delta_m,\qquad
 \Delta_{93}=\Delta_{95}=0.48,\qquad
 \Delta_E=1.9-0.43q_E.}
 \label{mawf:kappa}
\end{equation}
For E94/E95 the branch \(C_1=0.43\) is active at \(q_E\).
Retaining the original outer coefficients means selecting
\(\kappa_m=\sqrt{\sigma_{\eps,m}^{\rm orig}\Delta_m/q_m}\).
A prescribed common \(\kappa\) instead requires the single outer
replacement \(\sigma_{\eps,m}^{\rm eq}=\kappa^2q_m/\Delta_m\).
The exact values implied by the printed coefficients are listed in
\cref{mawf:outer-table}; they need not coincide with a nominal
\(\kappa\) used in the original empirical calibration.

\begin{table}[htbp]
\centering\small
\caption{Outer choices. The last column is the required one-coefficient
change when a common \(\kappa=0.41\) is selected.}
\label{mawf:outer-table}
\begin{tabular}{@{}lrrrr@{}}
\toprule
Model & \(q_m\) & original \(\sigma_\eps\)& compatible \(\kappa_m\)&
\(\sigma_\eps^{\,0.41}\)\\
\midrule
A93 &3.401265417&1.3&0.428323629&1.191151493\\
A95 &3.822000426&1.3&0.404060995&1.338496399\\
E94/E95 &3.324506456&1.2&0.412087472&1.187873345\\
\bottomrule
\end{tabular}
\end{table}

Use endpoint sampling and tail matching:
\begin{equation}
 \boxed{z_K=z_m=H,\qquad K=k_w(H)=k_o(H)>0.}
 \label{mawf:sampling}
\end{equation}
The subscripts \(w,o\) denote wall layer and outer region; subscript
\(0\) denotes their steady flat reference. At each positive current
sample \(K\), define counterfactual zero-forcing scales
\begin{equation}
 u_m(K)=\sqrt{\frac K{q_m}},\qquad
 t_0(K)=\frac K{q_m},\qquad
 A_{\eps,m}(K)=\frac{u_m(K)^3}{\kappa}.
 \label{mawf:scales}
\end{equation}
At the chosen inlet reference \(K=q_m\ust^2\), \(u_m=\ust\).
Away from that reference, \(t_0\) is not imposed as the actual
wall traction. Model and sample arguments are suppressed below.

Match the wall dissipation-tail amplitude and viscosity slope:
\begin{equation}
 \frac{K^{3/2}}{c_{\ell,m}^w}=\frac{u_m^3}{\kappa},\qquad
 c_{\mu,m}^wc_{\ell,m}^w\sqrt K=\kappa u_m.
\end{equation}
The two compatible auxiliary coefficients and their product are
\begin{equation}
 \boxed{c_{\ell,m}^w=\kappa q_m^{3/2},\qquad
 c_{\mu,m}^w=q_m^{-2},\qquad
 \alpha_m=c_{\mu,m}^wc_{\ell,m}^w=\frac{\kappa}{\sqrt{q_m}}.}
 \label{mawf:wall-coefficients}
\end{equation}
Consequently, with \(a=\alpha_m\sqrt K\),
\begin{equation}
 \boxed{a=\kappa u_m,\qquad
 \frac{K^{3/2}}{c_{\ell,m}^w}=A_{\eps,m}=\frac{t_0^2}{a}.}
 \label{mawf:identities}
\end{equation}
These two matching requirements uniquely determine the two auxiliary
coefficients. Traction continuity by itself does not require
equal inner and outer viscosity slopes.
The auxiliary \(c_\mu^w\) is constant for the selected branch;
it is not the current gradient-dependent outer \(C_\mu\).

\begin{table}[htbp]
\centering\small
\caption{Compatible wall coefficients. Each row uses the same
\(\kappa\) in the outer equations and the wall model.}
\label{mawf:coefficient-table}
\begin{tabular}{@{}llrrrr@{}}
\toprule
Model&outer choice&\(\kappa\)&\(c_\mu^w\)&\(c_\ell^w\)&\(\alpha\)\\
\midrule
A93&original&0.428323629&0.086440835&2.686784401&0.232247887\\
A95&original&0.404060995&0.068457106&3.019138428&0.206681480\\
E94/E95&original&0.412087472&0.090478551&2.497929531&0.226009045\\
\midrule
A93&common \(\kappa\)&0.410000000&0.086440835&2.571844117&0.222312353\\
A95&common \(\kappa\)&0.410000000&0.068457106&3.063514593&0.209719344\\
E94/E95&common \(\kappa\)&0.410000000&0.090478551&2.485276011&0.224864173\\
\bottomrule
\end{tabular}
\end{table}

Retain the original roughness-intercept correlation. With sand-grain
roughness \(h_r\ge0\), define, for \(h_r>0\),
\begin{equation}
 h^*=\frac{h_r\sqrt K}{\nu},\quad x=\frac{h^*}{70},\quad
 m_r=\max(0.5-0.4x^{0.7},\,1-0.79x^{-0.28}),\quad
 b=10.7(1-x^{m_r});
 \label{mawf:roughness}
\end{equation}
for \(h_r=0\), set \(b=10.7\). The wall coefficients are
\begin{equation}
 \begin{gathered}
 z_v=\frac{\nu b}{\sqrt K},\qquad
 \nu_{t,A}(z)=a(z-z_v)_+,\qquad D_A(z)=\nu+\nu_{t,A}(z),\\
 D_{k,w}(z)=\nu+\frac{\nu_{t,A}(z)}{\sigma_k}=D_A(z).
 \end{gathered}
 \label{mawf:viscosity}
\end{equation}
Here \((r)_+=\max(r,0)\), and the last equality uses \(\sigma_k=1\)
for all three models. These coefficients sample endpoint \(K\);
local transported \(k_w(z)\) is allowed to vary. Substituting local
\(k_w\) into the kernel would define a different model. Likewise,
an interior sample or average cannot simply replace the endpoint
sample once the inner reference is nonconstant.

\subsection{Momentum, production, and the outer velocity anchor}
\label{mawf:momentum}
For a flat quasi-steady column, put
\begin{equation}
 L_h(z)=(z-h_r)_+,\quad
 J_0(z)=\int_0^z\frac{\dd s}{D_A(s)},\quad
 J_1(z)=\int_0^z\frac{L_h(s)}{D_A(s)}\,\dd s,\quad
 J_2(z)=\int_0^z\frac{\dd s}{D_A(s)^2}.
 \label{mawf:J}
\end{equation}
With prescribed constant tangential forcing \(F\), retain
\begin{equation}
 (D_AU_A')'=F\,\mathbf1_{\{z>h_r\}},\qquad
 U_A(0)=0,\quad U_A(H)=U_H.
 \label{mawf:momentum-bvp}
\end{equation}
Integration yields the original analytical form, now evaluated
with the compatible viscosity:
\begin{equation}
 \boxed{\begin{aligned}
 t_A&=\frac{U_H-FJ_1(H)}{J_0(H)},&
 U_A(z)&=t_AJ_0(z)+FJ_1(z),\\
 P_A(z)&=\nu_{t,A}(z)
       \left|\frac{t_A+FL_h(z)}{D_A(z)}\right|^2 .
 \end{aligned}}
 \label{mawf:forced}
\end{equation}
The upper traction is \(T_H=t_A+FL_h(H)\), generally different
from ground traction \(t_A\). The vector tangential form uses
componentwise reconstruction and the Euclidean square in \(P_A\).
The coordinates remain the physical Cartesian coordinates.

At the counterfactual zero-forcing reference choose
\begin{equation}
 U_{H,0}=t_0J_0(H),\qquad U_{w,0}(z)=t_0J_0(z),\qquad
 p_0(z)=t_0^2\frac{\nu_{t,A}(z)}{D_A(z)^2}.
 \label{mawf:reference-production}
\end{equation}
Its integrated production is exactly
\begin{equation}
 \boxed{\mathcal P_0=\int_0^H p_0\,\dd z
     =t_0^2[J_0(H)-\nu J_2(H)].}
 \label{mawf:production-budget}
\end{equation}
Here \(p_0\) is pointwise reference production and
\(\mathcal P_0/H\) is its height mean. Neither is the actual
production of a general current flow. The molecular part of mean
viscous dissipation is excluded from turbulent production; hence
the term \(-\nu J_2\) must be retained.

For \(z\ge H\), the compatible outer reference is
\begin{equation}
 \boxed{\begin{aligned}
 U_{o,0}(z)&=t_0J_0(H)+\frac{u_m}{\kappa}\ln\frac zH,\\
 k_{o,0}(z)&=K,\qquad
 \eps_{o,0}(z)=\frac{A_{\eps,m}}z,\qquad
 \nu_{t,o,0}(z)=az .
 \end{aligned}}
 \label{mawf:outer-reference}
\end{equation}
Its aerodynamic intercept is
\begin{equation}
 z_{0,m}=H\exp[-\kappa u_mJ_0(H)]=H\exp[-aJ_0(H)].
 \label{mawf:z0}
\end{equation}
Keeping the same \(h_r\) after changing coefficients generally
changes \(z_{0,m}\). A separately prescribed \(z_0\) requires
\begin{equation}
 J_0(H)=\frac{\ln(H/z_0)}{\kappa\ust}
 \label{mawf:z0-calibration}
\end{equation}
at the selected inlet reference. A physical roughness input calibrated
there is subsequently held fixed as the endpoint sample evolves.

\subsection{Budget-determined dissipation and admissibility}
\label{mawf:cutoff-section}
Use a positive continuous dependent sink and no prescribed redistribution:
\begin{equation}
 \boxed{E_w(z)=\frac{A_{\eps,m}}{\max(z,z_c)},\qquad
 C_k^w(z)=0,\qquad 0<z_c\le H.}
 \label{mawf:dissipation}
\end{equation}
Its integral and upper trace are
\begin{equation}
 \int_0^HE_w\,\dd z=A_{\eps,m}\left(1+\ln\frac H{z_c}\right),
 \qquad E_w(H)=\frac{A_{\eps,m}}H.
 \label{mawf:dissipation-integral}
\end{equation}
Zero ground \(k\) flux and the constant outer reference require
\(\int_0^H(E_w-p_0)\,\dd z=0\). Thus
\begin{equation}
 \boxed{\Lambda=\frac{\mathcal P_0}{A_{\eps,m}}
       =a[J_0(H)-\nu J_2(H)],\qquad
 z_c=H\exp(1-\Lambda).}
 \label{mawf:cutoff}
\end{equation}
The integral in \cref{mawf:dissipation-integral} decreases strictly
with \(z_c\), tends to infinity as \(z_c\downarrow0\), and equals
\(A_{\eps,m}\) at \(z_c=H\). Therefore
\begin{equation}
 \boxed{\text{a unique }0<z_c\le H\text{ exists}
        \quad\Longleftrightarrow\quad \Lambda\ge1.}
 \label{mawf:gate}
\end{equation}
If the gate fails, this sink family cannot satisfy both the budget
and the outer trace. Clipping \(z_c\) to \(H\) loses the budget.

The original molecular cutoff
\(z_{\eps,A}=2c_\ell^w\nu/\sqrt K\) is replaced by
\cref{mawf:cutoff}. The plateau is \(A_{\eps,m}/z_c\).
Imposing also \(E_w(0)=2\nu K/z_c^2\) would force the original
molecular cutoff again. The replacement cutoff and plateau are
explicit modifications of the dissipation law.
The retained finite-viscosity contribution also gives
\begin{equation}
 \boxed{z_c=e\,z_{0,m}\exp[a\nu J_2(H)]>e\,z_{0,m},}
 \label{mawf:cutoff-z0}
\end{equation}
and the mechanical-work budget
\begin{equation}
 t_0U_{H,0}
 =\nu t_0^2J_2(H)+\int_0^Hp_0\,\dd z
 =\nu t_0^2J_2(H)+\int_0^HE_w\,\dd z.
 \label{mawf:work}
\end{equation}

\subsection{Nonconstant inner kinetic energy and every weak test mode}
\label{mawf:inner-k}
Retain the actual inner transport equation
\begin{equation}
 \partial_tk_w+\bm U_w\cdot\nabla k_w
 =\nabla\cdot(D_{k,w}\nabla k_w)+P_A-E_w.
 \label{mawf:k-transport}
\end{equation}
The steady parallel reference satisfies
\begin{equation}
 (D_{k,w}k_{w,0}')'=E_w-p_0,\qquad
 D_{k,w}(0)k_{w,0}'(0)=0,\qquad k_{w,0}(H)=K.
 \label{mawf:k-bvp}
\end{equation}
Define
\begin{equation}
 Q_0(z)=\int_0^z[E_w(s)-p_0(s)]\,\dd s.
\end{equation}
Two integrations give the unique reference profile
\begin{equation}
 \boxed{
 D_{k,w}k_{w,0}'=Q_0,\qquad
 k_{w,0}(z)=K-\int_z^H\frac{Q_0(s)}{D_{k,w}(s)}\,\dd s.}
 \label{mawf:k-profile}
\end{equation}
The physical upward diffusive flux is \(-Q_0\).
By definition \(Q_0(0)=0\), and the budget gives \(Q_0(H)=0\).
Hence
\begin{equation}
 k_{w,0}(H)=K,\qquad
 k_{w,0}'(0)=k_{w,0}'(H)=0.
 \label{mawf:k-endpoints}
\end{equation}
The ground value is determined by \cref{mawf:k-profile};
an additional \(k_w(0)=0\) generally overdetermines this reference.
The upper zero flux follows from the cutoff budget rather than
being an additional independent boundary condition on a fixed-cutoff
inner problem.

For any sufficiently regular test function \(\phi\),
\begin{equation}
 \boxed{
 -\int_0^H D_{k,w}k_{w,0}'\phi'\,\dd z
 +\int_0^H(p_0-E_w)\phi\,\dd z=0.}
 \label{mawf:weak}
\end{equation}
Indeed, substitution reduces the left side to
\(-[Q_0\phi]_0^H=0\). The constant test function expresses the
budget; for nonconstant modes, ordinary diffusion cancels nonzero
source moments. This is the mechanism required in a spectral/\(hp\)
wall element.

The analytical profile specifies a reference and an initialization;
the evolving \(k_w\) is transported, not reset to this profile.
Any equivalent prescribed redistribution flux from an earlier
flux-form implementation is also removed. The physical flux
\(-D_{k,w}\nabla k_w\) supplies the reference transport.
The cutoff continues to use counterfactual \(p_0(K)\).
For actual production,
\begin{equation}
 \int_0^H(P_A-E_w)\,\dd z=\int_0^H(P_A-p_0)\,\dd z
 \label{mawf:actual-budget}
\end{equation}
need not vanish. Determining the cutoff from actual \(P_A\)
would define a different model.

\subsection{Closed rough-branch formulas and positivity}
\label{mawf:rough}
Assume \(z_v\le0\), and put
\begin{equation}
 d=\nu-az_v\ge\nu,\qquad D_A=d+az,\qquad
 \zeta=\frac{az}{d},\quad L=\frac{aH}{d},\quad
 \beta=\frac{\nu}{d}\in(0,1].
 \label{mawf:rough-variables}
\end{equation}
Then
\begin{align}
 J_0(z)&=\frac1a\ln\frac{d+az}{d},&
 J_2(z)&=\frac1a\left(\frac1d-\frac1{d+az}\right),\\
 p_0(z)&=t_0^2\frac{d+az-\nu}{(d+az)^2},&
 \int_0^zp_0\,\dd s
 &=A_{\eps,m}\left[\ln\frac{d+az}{d}
          -\nu\left(\frac1d-\frac1{d+az}\right)\right].
 \label{mawf:rough-integrals}
\end{align}
Consequently,
\begin{equation}
 \boxed{\begin{gathered}
 \Lambda=\ln(1+L)-\frac{\beta L}{1+L},\qquad
 z_{0,m}=\frac{dH}{aH+d},\\
 z_c=e\,z_{0,m}
 \exp\left[\nu\left(\frac1d-\frac1{aH+d}\right)\right].
 \end{gathered}}
 \label{mawf:rough-cutoff}
\end{equation}
At a prescribed \(z_0<H\), the required rough-branch intercept is
\begin{equation}
 d=\frac{aH z_0}{H-z_0}.
 \label{mawf:d-calibration}
\end{equation}
It must be attainable by the roughness correlation and the selected
wall-layer geometry, normally \(h_r<H\).
Equal viscosity slopes do not remove the finite intercept:
\(D_A(H)=d+aH\), while the outer high-Re viscosity is \(aH\).

Define
\begin{equation}
 R_c(z)=
 \begin{cases}
 z/z_c,&0\le z\le z_c,\\
 1+\ln(z/z_c),&z_c\le z\le H.
 \end{cases}
\end{equation}
The cumulative defect is
\begin{equation}
 \boxed{Q_0(z)=A_{\eps,m}
 \left[R_c(z)-\ln\frac{d+az}{d}
       +\nu\left(\frac1d-\frac1{d+az}\right)\right].}
 \label{mawf:rough-Q}
\end{equation}

\paragraph{Closed kinetic-energy primitive.}
For \(c=az_c/d=L e^{1-\Lambda}\), set
\begin{equation}
 \mathcal F_<(\zeta)=
 \frac{\zeta-\ln(1+\zeta)}c-\frac12\ln^2(1+\zeta)
 +\beta\left[\ln(1+\zeta)+\frac1{1+\zeta}-1\right],
\end{equation}
and
\begin{equation}
 \begin{split}
 \mathcal G(\zeta)
 &=\left[1+\ln\frac{\zeta}{c}+\beta\right]\ln(1+\zeta)
   +\operatorname{Li}_2(-\zeta)
   -\frac12\ln^2(1+\zeta)+\frac{\beta}{1+\zeta},\\
 \operatorname{Li}_2(-\zeta)
 &=-\int_0^\zeta\frac{\ln(1+s)}s\,\dd s\qquad(\zeta\ge0).
 \end{split}
\end{equation}
Join the primitives:
\begin{equation}
 \mathcal F(\zeta)=
 \begin{cases}
 \mathcal F_<(\zeta),&\zeta\le c,\\
 \mathcal F_<(c)+\mathcal G(\zeta)-\mathcal G(c),&\zeta\ge c.
 \end{cases}
\end{equation}
Since
\(\mathcal F'(\zeta)=Q_0/[A_{\eps,m}(1+\zeta)]\),
\begin{equation}
 \boxed{k_{w,0}(z)=K+\frac{A_{\eps,m}}a
  [\mathcal F(az/d)-\mathcal F(aH/d)],\qquad
  \frac{A_{\eps,m}}a=\frac{u_m^2}{\kappa^2}.}
 \label{mawf:k-closed}
\end{equation}
Definite quadrature of \cref{mawf:k-profile} is an equivalent
evaluation. Near \(H\), it can avoid cancellation of close
primitive values.

\paragraph{A sufficient condition for monotone decrease.}
For \(z\ge z_c>0\), direct simplification using
\(A_{\eps,m}=t_0^2/a\) gives
\begin{equation}
 E_w-p_0
 =\frac{A_{\eps,m}[(d+\nu)az+d^2]}{z(d+az)^2}>0.
 \label{mawf:tail-defect}
\end{equation}
Moreover,
\begin{equation}
 p_0'(z)=\frac{a t_0^2[2\nu-(d+az)]}{(d+az)^3}.
 \label{mawf:production-derivative}
\end{equation}

\begin{proposition}[Monotone positive rough reference]
\label{mawf:monotone-proposition}
If \(d\ge2\nu\), \(K>0\), and \(\Lambda\ge1\), then
\begin{equation}
 Q_0(z)<0,\qquad k_{w,0}'(z)<0,\qquad
 k_{w,0}(z)>K>0\quad(0<z<H),
\end{equation}
and \(k_{w,0}(0)>K\).
\end{proposition}
\begin{proof}
The production decreases strictly for \(z>0\).
On the dissipation plateau, \(E_w-p_0\) therefore increases
strictly, and its value at \(z_c\) is positive by
\cref{mawf:tail-defect}. Its zero total integral forces one
negative-to-positive crossing strictly inside \((0,z_c)\).
The cumulative defect is negative until it returns to zero at
\(H\), and \cref{mawf:k-profile} gives the result.
This argument includes \(z_c=H\): the positive endpoint
defect still forces a crossing below \(H\).
\end{proof}
The one-crossing argument works for any positive \(D_{k,w}\).
The following stronger positivity statement uses specifically
\(\sigma_k=1\), hence \(D_{k,w}=D_A\).

\begin{proposition}[Positive reference on the whole rough branch]
\label{mawf:positive-proposition}
Assume the flat measure, constant physical \(\nu>0\),
\(d\ge\nu\), \(K>0\), \(\Lambda\ge1\), \(D_{k,w}=D_A\),
and \(A_{\eps,m}=t_0^2/a\). Then
\begin{equation}
 \boxed{k_{w,0}(z)>K\quad(0\le z<H),\qquad k_{w,0}(H)=K.}
 \label{mawf:positive}
\end{equation}
When \(\nu\le d<2\nu\), the profile can initially increase
before decreasing to \(K\); positivity does not require
monotone decrease.
\end{proposition}
\begin{proof}
Use \(L,\beta\) from \cref{mawf:rough-variables}, and set
\begin{equation}
 \begin{gathered}
 G_\beta(x)=\ln(1+x)-\frac{\beta x}{1+x},\qquad
 c=L e^{1-G_\beta(L)},\\
 \Phi(x;c)=
 \begin{cases}
 x/c-G_\beta(x),&x\le c,\\
 1+\ln(x/c)-G_\beta(x),&x\ge c.
 \end{cases}
 \end{gathered}
\end{equation}
Here \(Q_0=A_{\eps,m}\Phi\), \(\Phi(0)=\Phi(L)=0\), and
\begin{equation}
 k_{w,0}(z)-K=\frac{A_{\eps,m}}a[F(az/d)-F(L)],
 \qquad F(x)=\int_0^x\frac{\Phi(t;c)}{1+t}\,\dd t.
 \label{mawf:positivity-F}
\end{equation}
On the tail,
\(\Phi'=[1+(1+\beta)x]/[x(1+x)^2]>0\), so \(\Phi<0\)
before its endpoint zero. On the plateau,
\[
 \Phi'=\frac1c-p_\beta(x),\qquad
 p_\beta(x)=\frac{1+x-\beta}{(1+x)^2},\qquad
 p_\beta'(x)=\frac{2\beta-1-x}{(1+x)^3}.
\]
Thus \(p_\beta\) decreases or has only one maximum.
The allowed signs of \(\Phi\) are negative throughout the
interior, or positive initially and then negative.
This also holds when \(c=L\), since \(\Phi'(L)>0\).
Consequently \(F\) has no interior minimum.

For fixed \(\beta\), write \(I_\beta(L)=F(L)\) while accounting
for \(c=c(L)\). Since
\[
 \frac{c'(L)}{c(L)}
  =\frac{1+(1+\beta)L}{L(1+L)^2}>0,
\]
and \(\Phi(L)=0\), differentiation gives
\begin{equation}
 I_\beta'(L)=-\frac{c'(L)}{c(L)}
 \left[\frac{c-\ln(1+c)}c+\ln\frac{1+L}{1+c}\right]<0.
\end{equation}
Let \(L_1\) be the unique \(G_\beta(L_1)=1\).
At that lower admissible endpoint \(c=L_1\), and
\(T=\ln(1+L_1)\) satisfies \(1<T<2\). Direct integration gives
\begin{equation}
 I_\beta(L_1)
 =(2-T)\left(\frac{2-T}{2}-\frac T{L_1}\right)<0.
\end{equation}
The sign follows from
\(\ln(1+L_1)>2L_1/(L_1+2)\); the derivative of this difference
is \(L_1^2/[(1+L_1)(L_1+2)^2]>0\).
Hence \(F(L)<0=F(0)\) for every admissible \(L\).
The shape argument now gives \(F(x)>F(L)\) for \(x<L\),
which proves \cref{mawf:positive}.
\end{proof}
The proof covers \(d=\nu,\Lambda=1\) algebraically.
Whether such an endpoint is attainable with the original
roughness correlation and \(h_r<H\) is a separate geometric
admissibility check. The stronger proof is specific to the
flat measure and \(D_{k,w}=D_A\); it is not automatically valid
for another scalar diffusivity or a metric-weighted column.

\paragraph{Why the cutoff gate is not a general positivity test.}
On the smooth branch \(h_r=0\), set \(b=10.7\),
\(x=az/\nu\), \(v=10.7\alpha_m\), and
\(M=aH/\nu\). Let \(y_4=146.409759514\ldots\) solve
\[
 \ln(1+y_4)-\frac{y_4}{1+y_4}=4,
\]
and choose \(M=v+y_4\). Then \(\Lambda=4\) and
\(c=az_c/\nu=M e^{-3}\). With
\[
 B(x)=
 \begin{cases}
 0,&x\le v,\\
 \ln(1+x-v)-(x-v)/(1+x-v),&x>v,
 \end{cases}
 \qquad
 \widehat R_c(x)=
 \begin{cases}x/c,&x\le c,\\1+\ln(x/c),&x\ge c,\end{cases}
\]
the exact ground value is
\begin{equation}
 \frac{k_{w,0}(0)}K
 =1-\frac{1}{\kappa^2q_m}
   \int_0^M\frac{\widehat R_c(x)-B(x)}{1+(x-v)_+}\,\dd x.
 \label{mawf:smooth-counterexample}
\end{equation}
At common \(\kappa=0.41\), this gives
\begin{center}
\begin{tabular}{@{}lrrr@{}}
\toprule
Model&A93&A95&E94/E95\\
\midrule
\(k_{w,0}(0)/K\)&\(-0.371380517\)&\(-0.080424984\)&\(-0.436121275\)\\
\bottomrule
\end{tabular}
\end{center}
The budget is admissible, but the stationary kinetic energy is
inadmissible. Thus outside the rough positivity theorem,
the computed profile needs a separate positivity check.

\subsection{Conservative matching and normal pressure balance}
\label{mawf:matching}
At the flat reference,
\begin{equation}
 U_{w,0}(0)=0,\qquad U_{w,0}(H)=U_{o,0}(H),\qquad
 D_A(H)U_{w,0}'(H)=\nu_{t,o,0}(H)U_{o,0}'(H)=t_0.
 \label{mawf:traction}
\end{equation}
Viscosities and gradients need not be continuous.
The inner Reynolds shear is
\(-t_0\nu_{t,A}/D_A\), while outer Reynolds shear is \(-t_0\);
the inner molecular contribution completes total-traction balance.

Use local transported \(k_w\) in the isotropic inner stress.
With \(\Smat_w=(\nabla\bm U_w+\nabla\bm U_w^T)/2\) and
hydrostatic-reduced pressure \(\pi=p/\rho+gz\), the inner
covariance and total stress are
\begin{equation}
 \R_w=\frac23k_w\Id-2\nu_{t,A}\Smat_w,\qquad
 \bm T_w=-\left(\pi_w+\frac23k_w\right)\Id+2D_A\Smat_w .
 \label{mawf:stress}
\end{equation}
The steady vertical balance and its solution are
\begin{equation}
 \left(\pi_{w,0}+\frac23k_{w,0}\right)'=0,\qquad
 \boxed{\pi_{w,0}(z)=\pi_{o,0}+R_{zz,o,0}-\frac23k_{w,0}(z).}
 \label{mawf:pressure}
\end{equation}
Here \(\pi_{o,0}\) and \(R_{zz,o,0}\) are constant in the
outer reference. Define
\begin{equation}
 F_{93}=\frac{q_{93}^2}{1000+q_{93}^3},\qquad
 B_{95}=
 \frac{q_{95}^2\sqrt{1-9/(2q_{95}^2)}}{1+3q_{95}^2}.
\end{equation}
The original outer stresses give
\begin{equation}
 \boxed{\frac{\pi_{w,0}(H)-\pi_{o,0}}K
 =\begin{cases}
 -17F_{93}/3=-0.0630736219,&\text{A93},\\
 -B_{95}=-0.2710910066,&\text{A95},\\
 0,&\text{E94/E95}.
 \end{cases}}
 \label{mawf:pressure-offset}
\end{equation}
Define \(\Delta\pi_H=\pi_{w,0}(H)-\pi_{o,0}\).
Then \(\pi_{w,0}(z)=\pi_{o,0}+\Delta\pi_H
-\frac23[k_{w,0}(z)-K]\).
Even E94/E95 has a varying inner reduced mechanical pressure,
despite its zero interface offset. A pressure variable that already
includes isotropic turbulent stress must be translated consistently.

The \(k\) value and diffusive flux match:
\begin{equation}
 k_w(H)=k_o(H)=K,\qquad
 D_{k,w}(H)k_{w,0}'(H)
 =\frac{\nu_{t,o,0}(H)}{\sigma_k}k_{o,0}'(H)=0.
 \label{mawf:k-interface}
\end{equation}
For the evolving coupled system these are continuity and conservation
conditions, not independent Dirichlet and Neumann prescriptions on
an isolated outer equation. The outer dissipation equation receives
only
\begin{equation}
 \eps_o(H)=E_w(H)=\frac{A_{\eps,m}}H .
 \label{mawf:eps-interface}
\end{equation}
There is no additional two-sided dissipation-flux condition,
because no inner dissipation equation is solved.
When \(z_c=H\), the inner dissipation derivative is zero and the
outer derivative is \(-A_{\eps,m}/H^2\); their traces still agree.

\subsection{Distinction from the constant-inner-\texorpdfstring{\(k\)}{k} EP-AWF}
\label{mawf:comparison}
At identical sampled \(K\), \(h_r,H,\nu\), compatible coefficients
and cutoff, the constant-inner-\(k\) construction uses
\[
 E_w=\frac{A_{\eps,m}}{\max(z,z_c)},\qquad
 C_k^w=E_w-p_0,\qquad S_{k,w}^{\rm EP}=P_A-p_0.
\]
The modified construction uses
\[
 E_w=\frac{A_{\eps,m}}{\max(z,z_c)},\qquad
 C_k^w=0,\qquad S_{k,w}=P_A-E_w.
\]
Both sources have the same layer integral, but their local
balances differ. At the reference, the first uses prescribed
redistribution to support constant \(k=K\); the second uses
ordinary diffusion of \cref{mawf:k-profile}.
The physical reference transport \(-Q_0\) is the same,
while the evolving models are different.

Setting the host's old switch to \(\chi=0\) does not select this
variant when that switch also restores the original-cutoff
dissipation. The present variant retains the budget cutoff
while removing both the prescribed redistribution source and
its equivalent auxiliary flux. Its scalar policy has no
spatial \(\chi\) fade. The earlier bounded-redistribution
positivity remedy would be a further model change if retained.

\subsection{A specified extension to curved terrain}
\label{mawf:curved}
Use the conservative curved momentum completion with the new
endpoint sample \(K=k|_{\Gamma_H}\), compatible coefficients,
metric-weighted reference budget and \(C_k^w=0\).
The equations are restated here so that the scalar replacement
and roughness-force prescription are explicit.

Let \(\bm n\) point from the ground into the fluid and use the
normal map
\begin{equation}
 \bm x(\bm\xi,r)=\bm X(\bm\xi)+r\bm n(\bm\xi),\qquad
 P=\Id-\bm n\otimes\bm n,\quad B_g=\nabla_\Gamma\bm n,\quad
 A_g=I_T+rB_g,\quad J=\det_TA_g.
 \label{mawf:map}
\end{equation}
The inverses of \(A_g\) act tangentially and are extended by zero
normally. The principal stretches \(1+r\gamma_i\) must be
positive, and the normal map must be globally one-to-one on
\(0\le r\le H(\bm\xi)\). The volume measure is
\(J\,\dd S_\Gamma\,\dd r\). Surface derivatives below are
taken at fixed \(r\) before restriction to the interface.

\paragraph{Reference cutoff in the volume measure.}
Set
\begin{equation}
 E_w(r)=\frac{A_{\eps,m}(K)}{\max(r,r_c)},\qquad
 p_0(r;K)=t_0(K)^2\frac{\nu_{t,A}(r;K)}{D_A(r;K)^2}.
\end{equation}
Use counterfactual zero-forcing production in
\begin{equation}
 \boxed{\int_0^H J(r)E_w(r)\,\dd r
        =\mathcal P_J:=\int_0^H J(r)p_0(r;K)\,\dd r.}
 \label{mawf:metric-budget}
\end{equation}
For \(\mathcal B(c)=A_{\eps,m}\int_0^H J/\max(r,c)\,\dd r\),
\[
 \mathcal B'(c)=-\frac{A_{\eps,m}}{c^2}\int_0^cJ\,\dd r<0,
 \qquad \mathcal B(c)\longrightarrow+\infty\quad(c\downarrow0).
\]
Hence a unique \(0<r_c\le H\) exists exactly when
\begin{equation}
 \boxed{\mathcal P_J\ge
        \frac{A_{\eps,m}}H\int_0^H J\,\dd r.}
 \label{mawf:metric-gate}
\end{equation}
At a fixed footpoint, \(J=(1+\gamma_1r)(1+\gamma_2r)\), and
\begin{equation}
 \frac{\mathcal B(c)}{A_{\eps,m}}
 =1+\ln\frac Hc
 +(\gamma_1+\gamma_2)\left(H-\frac c2\right)
 +\gamma_1\gamma_2\left(\frac{H^2}{2}-\frac{c^2}{6}\right).
 \label{mawf:metric-explicit}
\end{equation}
One monotone scalar root solve replaces the flat exponential.

\paragraph{Actual continuity, momentum, forcing and scalar transport.}
Write \(\bm U_w=\bm u_t+w\bm n\) and
\(\bm v=A_g^{-1}\bm u_t\). With Cartesian velocity components,
\begin{equation}
 \partial_r(Jw)+\operatorname{div}_\Gamma(J\bm v)=0,
 \label{mawf:continuity}
\end{equation}
\begin{equation}
 \begin{split}
 J\partial_t\bm U_w
 &+\partial_r\{J(w\bm U_w-T_w\bm n)\}\\
 &+\operatorname{Div}_\Gamma
   \{J(\bm U_w\otimes\bm v-T_wA_g^{-T})\}=J\bm f_w .
 \end{split}
 \label{mawf:curved-momentum}
\end{equation}
Here \(T_w\) is \cref{mawf:stress};
the outer stress is \(T_o=-\pi_o\Id-\R_m\).
The operator \(\operatorname{Div}_\Gamma\) takes the surface
divergence of each Cartesian row. These fluxes already contain
the geometric effects of stress and pressure.

Retain the explicit roughness-force completion
\begin{equation}
 \begin{split}
 a_r(r)&=\mathbf1_{\{r>h_r\}},\qquad
 \bm f_w=\bm f_\pi+(1-a_r)\bm F_A,\\
 \bm F_A&=\mathcal E_H^+P
 \left[\operatorname{Div}_r(\bm u_{o,t}\otimes\bm u_{o,t})
             +\nabla_r\pi_o-\bm f_\pi\right],
 \qquad \bm u_{o,t}=P\bm U_o .
 \end{split}
 \label{mawf:rough-force}
\end{equation}
Here \(\bm f_\pi\) is the remaining body acceleration after the
same hydrostatic-pressure convention used outside,
\(\nabla_rf=A_g^{-T}\nabla_\Gamma f\), and
\[
 \operatorname{Div}_r Q
 =J^{-1}\operatorname{Div}_\Gamma(JQA_g^{-T})
 \quad\text{for }Q\bm n=0.
\]
The sampler \(\mathcal E_H^+\) takes one-sided outer values after
derivatives at fixed \(r\) have been evaluated; extend the resulting
\(\bm F_A\) constantly along each normal ray.
Its work has no guaranteed drag sign and must be retained in
energy accounting.

The flat scalar reconstruction is the declared quasi-steady
parallel-flow reduction of this completion: inner storage,
normal advection, lateral viscous transport and the relevant
\(\nabla_\parallel w\) contributions are neglected, and the
remaining convection/pressure/body-force balance is replaced
by the height-constant sample \(\bm F_A\).
This also replaces the inner effective-pressure gradient by the
sampled outer mechanical-pressure gradient; the discarded
\(\nabla_\parallel(\pi_w+2k_w/3-\pi_o)\) is generally nonzero.
Under those approximations
\((D_A\bm U_\parallel')'=a_r\bm F_A\).
For the curved or unsteady patch, use the full equation
\cref{mawf:curved-momentum}, without additionally imposing that
scalar reduction. All discarded terms vanish on the homogeneous
zero-forcing reference.

The modified actual inner \(k\) equation is
\begin{equation}
 \boxed{\begin{aligned}
 J\partial_t k_w
 &+\partial_r\{J(wk_w-D_{k,w}\partial_rk_w)\}\\
 &+\operatorname{div}_\Gamma
 \{J[\bm v k_w-D_{k,w}A_g^{-1}A_g^{-T}\nabla_\Gamma k_w]\}
 =J(P_A-E_w).
 \end{aligned}}
 \label{mawf:curved-k}
\end{equation}
For incompressible flow,
\[
 P_A=2\nu_{t,A}\Smat_w:\Smat_w .
\]
The isotropic stress uses local \(k_w\), while the viscosity
uses sampled \(K\). If incompressibility is enforced only weakly,
the production contraction \(-\R_w:\nabla\bm U_w\) also contains
\(-(2k_w/3)\nabla\cdot\bm U_w\); its treatment must agree with
the discrete stress-work and pressure convention.
The base-normal tangential traction is
\begin{equation}
 \bm t_A=D_A\left[
 \partial_r\bm u_t+A_g^{-T}\nabla_\Gamma w-B_gA_g^{-1}\bm u_t
 \right].
 \label{mawf:curved-traction}
\end{equation}
Curvature, normal motion and normal pressure balance are thus
retained in the full patch.

The actual source budget remains
\[
 \int_0^HJ(P_A-E_w)\,\dd r
 =\int_0^HJ(P_A-p_0)\,\dd r,
\]
which need not vanish. A frozen reference column can be defined by
\begin{equation}
 Q_J(r)=\int_0^rJ(s)[E_w(s)-p_0(s)]\,\dd s,\qquad
 k_{\rm col}(r)=K-\int_r^H
        \frac{Q_J(s)}{J(s)D_{k,w}(s)}\,\dd s.
 \label{mawf:curved-column}
\end{equation}
Its base-normal endpoint fluxes vanish.
The one-crossing proof for \(d\ge2\nu\) extends to this column
for positive \(J,D_{k,w}\).
The stronger whole-rough-branch proof is not automatically
metric weighted. The column balances only the reference source
and normal diffusion; it is not generally an equilibrium of
the complete terrain-flow equations.

\paragraph{Actual interface geometry and shared fluxes.}
For variable \(H\), put
\begin{equation}
 \bm q_H=A_{g,H}^{-T}\nabla_\Gamma H,\qquad
 \beta_H=\sqrt{1+|\bm q_H|^2},\qquad
 \bm N_H=\frac{\bm n-\bm q_H}{\beta_H}.
 \label{mawf:interface-normal}
\end{equation}
Use common velocity and total traction,
\(\bm U_w=\bm U_o=\widehat{\bm U}\) and
\(T_w\bm N_H=T_o\bm N_H\), with one conservative momentum
flux and opposite orientations on the two sides.
The mass flux per unit ground area is
\(m_H=J_H(w-\bm u_t\cdot\bm q_H)\).
At a common \(k\) trace \(\widehat k\), match
\begin{equation}
 J_H\beta_H(\bm U k-D_k\nabla k)\cdot\bm N_H
 =m_H\widehat k-J_HD_k
 \left[\partial_rk-\bm q_H\cdot A_{g,H}^{-T}\nabla_\Gamma k\right].
 \label{mawf:curved-k-flux}
\end{equation}
Zero base-normal column flux alone does not establish zero flux
through a tilted interface. Surface differentiation of the trace
obeys
\(\nabla_\Gamma K=(\nabla_\Gamma k_w)|_H+
(\partial_rk_w)|_H\nabla_\Gamma H\).
The outer dissipation trace is \(E_w(H)=A_{\eps,m}/H\).
Full normal-traction matching replaces the flat pressure offset;
there is still no independent inner dissipation-flux prescription.

\subsection{Reference realizability, positivity, and stability limits}
\label{mawf:stability}
On a positive rough reference the smallest inner covariance
eigenvalue is
\begin{equation}
 \lambda_{\min}(\R_{w,0})
 =\frac23k_{w,0}-t_0\frac{\nu_{t,A}}{D_A}
 >K\left(\frac23-\frac1{q_m}\right)>0.
 \label{mawf:covariance}
\end{equation}
All three \(q_m\) exceed \(3/2\). Their original outer reference
covariances are also positive definite. These are reference-state
results, not global guarantees for forced or separated states.

With endpoint \(K\), coefficients, cutoff and reference production
held fixed, a wall-normal perturbation \(v=k_w-k_{w,0}\) obeys
\[
 v_t=(D_{k,w}v')',\qquad v(H)=0,\qquad D_{k,w}v'(0)=0,
\]
and
\begin{equation}
 \frac12\frac{\dd}{\dd t}\int_0^H v^2\,\dd z
 =-\int_0^H D_{k,w}(v')^2\,\dd z\le0.
 \label{mawf:scalar-stability}
\end{equation}
This conditional scalar estimate is separate from full coupled
stability.

Dynamic nonnegativity of actual \(k_w\) is not automatic:
at a flat interior zero minimum, positive endpoint \(K\), and
locally zero production and diffusion,
\[
 \partial_t k_w=-E_w(K)<0.
\]
Thus the positive rough reference does not make the evolving
sampled-\(K\) sink positivity preserving. An additional source
modification or numerical positivity device must be identified
and checked separately if introduced.

The wall changes do not alter the original outer constitutive
domains or incremental momentum response. A93 can lose covariance
realizability away from the reference, A95 retains its real-radicand
restriction, and positive covariance in E94/E95 does not ensure
positive incremental momentum diffusion. The previously established
outer constitutive restrictions therefore remain relevant.

\subsection{Illustrative references and implementation implications}
\label{mawf:verification}
Take
\begin{equation}
 \ust=0.4\ \mathrm{m\,s^{-1}},\quad
 \nu=1.5\times10^{-5}\ \mathrm{m^2\,s^{-1}},\quad
 h_r=0.1\ \mathrm m,\quad H=2\ \mathrm m,\quad z_K=z_m=H.
 \label{mawf:inputs}
\end{equation}
Using the original roughness correlation gives
\cref{mawf:numerical-table}. The same sand-grain input is retained,
so aerodynamic roughness is an output rather than a common
prescribed input.

\begin{table}[htbp]
\centering\small
\caption{Analytical rough references. \(K,k(0)\) are in
\(\mathrm{m^2\,s^{-2}}\), and lengths are in metres.}
\label{mawf:numerical-table}
\begin{tabular}{@{}llrrrrr@{}}
\toprule
Model&\(\kappa\) choice&\(K\)&\(k(0)\)&\(\Lambda\)&\(z_c\)&\(z_{0,m}\)\\
\midrule
A93&original&0.54420247&1.08630845&5.90564847&0.01480928&0.00536024\\
A95&original&0.61152007&1.22075495&5.89870963&0.01491240&0.00539286\\
E94/E95&original&0.53192103&1.11770080&5.90535573&0.01481362&0.00535834\\
\midrule
A93&common \(0.41\)&0.54420247&1.13595953&5.90420900&0.01483061&0.00536413\\
A95&common \(0.41\)&0.61152007&1.20319430&5.89920134&0.01490507&0.00539152\\
E94/E95&common \(0.41\)&0.53192103&1.12369422&5.90518519&0.01481614&0.00535880\\
\bottomrule
\end{tabular}
\end{table}

All six cases satisfy \(d\ge2\nu\) and \(\Lambda>1\), so their
positive monotone profiles have an analytical guarantee.
The interface pressure offsets are
\(-0.034324820606\), \(-0.165777590815\), and \(0\)
\(\mathrm{m^2\,s^{-2}}\), respectively.
At fixed \(\ust\) these offsets do not change with \(\kappa\),
because \(K=q_m\ust^2\) is unchanged; the inner pressure profile
does change with \(k_w(z)\).

Independent quadrature of the production and dissipation budgets,
the cumulative defect, and the closed primitive agrees with the
analytical expressions. Separately integrated source and diffusion
moments for Legendre degrees \(0,\ldots,6\) give a maximum absolute
weak residual of \(2.33\times10^{-16}\) across these cases;
closed \(k\) versus quadrature differs by at most
\(6.89\times10^{-15}\) in SI units. These are checks of the
formulas, not CFD or terrain validation.

The model-specific \(\kappa\) affects the viscosity slope,
dissipation amplitude, cutoff and inner profile. A further useful
identity is
\[
 \frac{\kappa^2[k_{w,0}(z)-K]}{u_m^2}
 =\mathcal F(az/d)-\mathcal F(aH/d).
\]
At common \(\kappa,\ust,H,z_0,\nu\), if each roughness input
is calibrated to the same \(z_0\), \cref{mawf:d-calibration}
makes \(a,d\) identical. Then the zero-forcing wall velocity,
production, sink and cumulative defect coincide across the models,
and their \(k\) profiles differ only by \(K=q_m\ust^2\).
Their forced responses can still differ through the calibrated
\(h_r\) in the forcing cutoff.

For a spectral/\(hp\) implementation, the exact continuum weak
identity \cref{mawf:weak} does not by itself make a finite polynomial
interpolant stationary. The ordinary diffusive flux, nonzero source
moments and shared interface terms must be represented and integrated
consistently. A local discrete wall solve using the same source,
quadrature and boundary contract can supply a stationary discrete
profile. Checking only the cell-mean energy budget is insufficient.

The selected variant therefore consists of endpoint sampling,
the two compatible coefficients, the budget cutoff and its plateau,
retained finite-\(\nu\) transport with \(P_A-E_w\), and conservative
velocity/stress/scalar coupling. Its exactness claim is the
flat neutral reference of these specified equations. Positivity
outside the rough theorem, global coupled stability and fully
discrete preservation remain separate requirements.

% The main note already calls \appendix; do not call it again here.
% Uses the main note's amsmath/amssymb/amsthm/bm/booktabs/natbib/cleveref
% packages, proposition environment, and \eps, \ust, \dd, \Id, \R, \Smat.
% Bibliography keys: Shih1993, Shih1995ASM, Shih1994EVM, Shih1995EVM,
% and SCI2006. All labels introduced here have the unique prefix mawf:.
% No document class, document environment, or bibliography is added.

\section{AWF with coefficients compatible with A93, A95 and E94/E95,
a budget-determined cutoff, and a nonconstant inner kinetic-energy profile}
\label{mawf:appendix}

This appendix specifies a modified analytical wall function for each of
the three outer closures A93, A95 and E94/E95
\citep{Shih1993,Shih1995ASM,Shih1994EVM,Shih1995EVM}.
The objective is to retain its exact neutral outer equilibrium while
allowing kinetic energy to vary inside the modeled wall layer.
Two auxiliary coefficients are matched to the selected outer branch;
the dissipation cutoff is determined by the reference energy budget;
and the prescribed redistribution source is removed. Ordinary diffusion
of the transported inner kinetic energy then balances the local
production--dissipation difference.

The rough-wall viscosity kernel and analytical momentum reconstruction
start from the AWF of Suga, Craft and Iacovides \citep{SCI2006}.
The compatible coefficients, endpoint sampling, budget cutoff and
variable inner reference specified below are explicit modifications.
Selecting this variant replaces the earlier wall-coefficient,
sampling and \(\chi\)-dependent scalar-source choices. Its own
dissipation law is retained with \(C_k^w=0\).

\subsection{Scope, endpoint sampling, and compatible coefficients}
\label{mawf:coefficients}

Use the selected high-Re outer equations: molecular momentum and scalar
diffusion are omitted from the outer analytical reference, and the
E94/E95 dissipation denominator is reduced from
\(k+\sqrt{\nu\eps}\) to \(k\) in that outer limit.
The inner AWF retains the fixed physical \(\nu>0\) in its momentum,
production and scalar-diffusion formulas. These are two declared
equation regimes; the result is not an exact logarithmic solution of
an unchanged finite-Re outer system.

Only inner \(k_w\) is transported. Inner dissipation is a dependent
total sink, and no inner dissipation transport equation is introduced.
The inner stress uses the auxiliary linear viscosity specified below;
the original A93/A95 nonlinear stress laws remain the outer closures.
Coordinates are \(x\) streamwise, \(y\) spanwise, and \(z\) vertical.
All stresses are per unit density.

Let \(m\in\{93,95,E\}\) identify the outer model.
Its neutral equilibrium ratio \(q_m=k_{o,0}/\ust^2\) is
\begin{equation}
 \boxed{
 q_{93}=\frac{57+\sqrt{6249}}{40},\qquad
 q_{95}=\frac{3+\sqrt{61}}{2\sqrt2},\qquad
 q_E=\frac{3+\sqrt{41}}{2\sqrt2}.}
 \label{mawf:q}
\end{equation}
These roots follow from \(C_{\mu,m}(q_m)q_m^2=1\) in simple shear.
The nonlinear corrections of A93 and A95 have zero \(xz\) component
there, but the full Reynolds shear is \(R_{xz,o,0}=-\ust^2\).

The outer dissipation equation additionally requires
\begin{equation}
 \boxed{\kappa^2q_m=\sigma_{\eps,m}\Delta_m,\qquad
 \Delta_{93}=\Delta_{95}=0.48,\qquad
 \Delta_E=1.9-0.43q_E.}
 \label{mawf:kappa}
\end{equation}
For E94/E95 the branch \(C_1=0.43\) is active at \(q_E\).
Retaining the original outer coefficients means selecting
\(\kappa_m=\sqrt{\sigma_{\eps,m}^{\rm orig}\Delta_m/q_m}\).
A prescribed common \(\kappa\) instead requires the single outer
replacement \(\sigma_{\eps,m}^{\rm eq}=\kappa^2q_m/\Delta_m\).
The exact values implied by the printed coefficients are listed in
\cref{mawf:outer-table}; they need not coincide with a nominal
\(\kappa\) used in the original empirical calibration.

\begin{table}[htbp]
\centering\small
\caption{Outer choices. The last column is the required one-coefficient
change when a common \(\kappa=0.41\) is selected.}
\label{mawf:outer-table}
\begin{tabular}{@{}lrrrr@{}}
\toprule
Model & \(q_m\) & original \(\sigma_\eps\)& compatible \(\kappa_m\)&
\(\sigma_\eps^{\,0.41}\)\\
\midrule
A93 &3.401265417&1.3&0.428323629&1.191151493\\
A95 &3.822000426&1.3&0.404060995&1.338496399\\
E94/E95 &3.324506456&1.2&0.412087472&1.187873345\\
\bottomrule
\end{tabular}
\end{table}

Use endpoint sampling and tail matching:
\begin{equation}
 \boxed{z_K=z_m=H,\qquad K=k_w(H)=k_o(H)>0.}
 \label{mawf:sampling}
\end{equation}
The subscripts \(w,o\) denote wall layer and outer region; subscript
\(0\) denotes their steady flat reference. At each positive current
sample \(K\), define counterfactual zero-forcing scales
\begin{equation}
 u_m(K)=\sqrt{\frac K{q_m}},\qquad
 t_0(K)=\frac K{q_m},\qquad
 A_{\eps,m}(K)=\frac{u_m(K)^3}{\kappa}.
 \label{mawf:scales}
\end{equation}
At the chosen inlet reference \(K=q_m\ust^2\), \(u_m=\ust\).
Away from that reference, \(t_0\) is not imposed as the actual
wall traction. Model and sample arguments are suppressed below.

Match the wall dissipation-tail amplitude and viscosity slope:
\begin{equation}
 \frac{K^{3/2}}{c_{\ell,m}^w}=\frac{u_m^3}{\kappa},\qquad
 c_{\mu,m}^wc_{\ell,m}^w\sqrt K=\kappa u_m.
\end{equation}
The two compatible auxiliary coefficients and their product are
\begin{equation}
 \boxed{c_{\ell,m}^w=\kappa q_m^{3/2},\qquad
 c_{\mu,m}^w=q_m^{-2},\qquad
 \alpha_m=c_{\mu,m}^wc_{\ell,m}^w=\frac{\kappa}{\sqrt{q_m}}.}
 \label{mawf:wall-coefficients}
\end{equation}
Consequently, with \(a=\alpha_m\sqrt K\),
\begin{equation}
 \boxed{a=\kappa u_m,\qquad
 \frac{K^{3/2}}{c_{\ell,m}^w}=A_{\eps,m}=\frac{t_0^2}{a}.}
 \label{mawf:identities}
\end{equation}
These two matching requirements uniquely determine the two auxiliary
coefficients. Traction continuity by itself does not require
equal inner and outer viscosity slopes.
The auxiliary \(c_\mu^w\) is constant for the selected branch;
it is not the current gradient-dependent outer \(C_\mu\).

\begin{table}[htbp]
\centering\small
\caption{Compatible wall coefficients. Each row uses the same
\(\kappa\) in the outer equations and the wall model.}
\label{mawf:coefficient-table}
\begin{tabular}{@{}llrrrr@{}}
\toprule
Model&outer choice&\(\kappa\)&\(c_\mu^w\)&\(c_\ell^w\)&\(\alpha\)\\
\midrule
A93&original&0.428323629&0.086440835&2.686784401&0.232247887\\
A95&original&0.404060995&0.068457106&3.019138428&0.206681480\\
E94/E95&original&0.412087472&0.090478551&2.497929531&0.226009045\\
\midrule
A93&common \(\kappa\)&0.410000000&0.086440835&2.571844117&0.222312353\\
A95&common \(\kappa\)&0.410000000&0.068457106&3.063514593&0.209719344\\
E94/E95&common \(\kappa\)&0.410000000&0.090478551&2.485276011&0.224864173\\
\bottomrule
\end{tabular}
\end{table}

Retain the original roughness-intercept correlation. With sand-grain
roughness \(h_r\ge0\), define, for \(h_r>0\),
\begin{equation}
 h^*=\frac{h_r\sqrt K}{\nu},\quad x=\frac{h^*}{70},\quad
 m_r=\max(0.5-0.4x^{0.7},\,1-0.79x^{-0.28}),\quad
 b=10.7(1-x^{m_r});
 \label{mawf:roughness}
\end{equation}
for \(h_r=0\), set \(b=10.7\). The wall coefficients are
\begin{equation}
 \begin{gathered}
 z_v=\frac{\nu b}{\sqrt K},\qquad
 \nu_{t,A}(z)=a(z-z_v)_+,\qquad D_A(z)=\nu+\nu_{t,A}(z),\\
 D_{k,w}(z)=\nu+\frac{\nu_{t,A}(z)}{\sigma_k}=D_A(z).
 \end{gathered}
 \label{mawf:viscosity}
\end{equation}
Here \((r)_+=\max(r,0)\), and the last equality uses \(\sigma_k=1\)
for all three models. These coefficients sample endpoint \(K\);
local transported \(k_w(z)\) is allowed to vary. Substituting local
\(k_w\) into the kernel would define a different model. Likewise,
an interior sample or average cannot simply replace the endpoint
sample once the inner reference is nonconstant.

\subsection{Momentum, production, and the outer velocity anchor}
\label{mawf:momentum}
For a flat quasi-steady column, put
\begin{equation}
 L_h(z)=(z-h_r)_+,\quad
 J_0(z)=\int_0^z\frac{\dd s}{D_A(s)},\quad
 J_1(z)=\int_0^z\frac{L_h(s)}{D_A(s)}\,\dd s,\quad
 J_2(z)=\int_0^z\frac{\dd s}{D_A(s)^2}.
 \label{mawf:J}
\end{equation}
With prescribed constant tangential forcing \(F\), retain
\begin{equation}
 (D_AU_A')'=F\,\mathbf1_{\{z>h_r\}},\qquad
 U_A(0)=0,\quad U_A(H)=U_H.
 \label{mawf:momentum-bvp}
\end{equation}
Integration yields the original analytical form, now evaluated
with the compatible viscosity:
\begin{equation}
 \boxed{\begin{aligned}
 t_A&=\frac{U_H-FJ_1(H)}{J_0(H)},&
 U_A(z)&=t_AJ_0(z)+FJ_1(z),\\
 P_A(z)&=\nu_{t,A}(z)
       \left|\frac{t_A+FL_h(z)}{D_A(z)}\right|^2 .
 \end{aligned}}
 \label{mawf:forced}
\end{equation}
The upper traction is \(T_H=t_A+FL_h(H)\), generally different
from ground traction \(t_A\). The vector tangential form uses
componentwise reconstruction and the Euclidean square in \(P_A\).
The coordinates remain the physical Cartesian coordinates.

At the counterfactual zero-forcing reference choose
\begin{equation}
 U_{H,0}=t_0J_0(H),\qquad U_{w,0}(z)=t_0J_0(z),\qquad
 p_0(z)=t_0^2\frac{\nu_{t,A}(z)}{D_A(z)^2}.
 \label{mawf:reference-production}
\end{equation}
Its integrated production is exactly
\begin{equation}
 \boxed{\mathcal P_0=\int_0^H p_0\,\dd z
     =t_0^2[J_0(H)-\nu J_2(H)].}
 \label{mawf:production-budget}
\end{equation}
Here \(p_0\) is pointwise reference production and
\(\mathcal P_0/H\) is its height mean. Neither is the actual
production of a general current flow. The molecular part of mean
viscous dissipation is excluded from turbulent production; hence
the term \(-\nu J_2\) must be retained.

For \(z\ge H\), the compatible outer reference is
\begin{equation}
 \boxed{\begin{aligned}
 U_{o,0}(z)&=t_0J_0(H)+\frac{u_m}{\kappa}\ln\frac zH,\\
 k_{o,0}(z)&=K,\qquad
 \eps_{o,0}(z)=\frac{A_{\eps,m}}z,\qquad
 \nu_{t,o,0}(z)=az .
 \end{aligned}}
 \label{mawf:outer-reference}
\end{equation}
Its aerodynamic intercept is
\begin{equation}
 z_{0,m}=H\exp[-\kappa u_mJ_0(H)]=H\exp[-aJ_0(H)].
 \label{mawf:z0}
\end{equation}
Keeping the same \(h_r\) after changing coefficients generally
changes \(z_{0,m}\). A separately prescribed \(z_0\) requires
\begin{equation}
 J_0(H)=\frac{\ln(H/z_0)}{\kappa\ust}
 \label{mawf:z0-calibration}
\end{equation}
at the selected inlet reference. A physical roughness input calibrated
there is subsequently held fixed as the endpoint sample evolves.

\subsection{Budget-determined dissipation and admissibility}
\label{mawf:cutoff-section}
Use a positive continuous dependent sink and no prescribed redistribution:
\begin{equation}
 \boxed{E_w(z)=\frac{A_{\eps,m}}{\max(z,z_c)},\qquad
 C_k^w(z)=0,\qquad 0<z_c\le H.}
 \label{mawf:dissipation}
\end{equation}
Its integral and upper trace are
\begin{equation}
 \int_0^HE_w\,\dd z=A_{\eps,m}\left(1+\ln\frac H{z_c}\right),
 \qquad E_w(H)=\frac{A_{\eps,m}}H.
 \label{mawf:dissipation-integral}
\end{equation}
Zero ground \(k\) flux and the constant outer reference require
\(\int_0^H(E_w-p_0)\,\dd z=0\). Thus
\begin{equation}
 \boxed{\Lambda=\frac{\mathcal P_0}{A_{\eps,m}}
       =a[J_0(H)-\nu J_2(H)],\qquad
 z_c=H\exp(1-\Lambda).}
 \label{mawf:cutoff}
\end{equation}
The integral in \cref{mawf:dissipation-integral} decreases strictly
with \(z_c\), tends to infinity as \(z_c\downarrow0\), and equals
\(A_{\eps,m}\) at \(z_c=H\). Therefore
\begin{equation}
 \boxed{\text{a unique }0<z_c\le H\text{ exists}
        \quad\Longleftrightarrow\quad \Lambda\ge1.}
 \label{mawf:gate}
\end{equation}
If the gate fails, this sink family cannot satisfy both the budget
and the outer trace. Clipping \(z_c\) to \(H\) loses the budget.

The original molecular cutoff
\(z_{\eps,A}=2c_\ell^w\nu/\sqrt K\) is replaced by
\cref{mawf:cutoff}. The plateau is \(A_{\eps,m}/z_c\).
Imposing also \(E_w(0)=2\nu K/z_c^2\) would force the original
molecular cutoff again. The replacement cutoff and plateau are
explicit modifications of the dissipation law.
The retained finite-viscosity contribution also gives
\begin{equation}
 \boxed{z_c=e\,z_{0,m}\exp[a\nu J_2(H)]>e\,z_{0,m},}
 \label{mawf:cutoff-z0}
\end{equation}
and the mechanical-work budget
\begin{equation}
 t_0U_{H,0}
 =\nu t_0^2J_2(H)+\int_0^Hp_0\,\dd z
 =\nu t_0^2J_2(H)+\int_0^HE_w\,\dd z.
 \label{mawf:work}
\end{equation}

\subsection{Nonconstant inner kinetic energy and every weak test mode}
\label{mawf:inner-k}
Retain the actual inner transport equation
\begin{equation}
 \partial_tk_w+\bm U_w\cdot\nabla k_w
 =\nabla\cdot(D_{k,w}\nabla k_w)+P_A-E_w.
 \label{mawf:k-transport}
\end{equation}
The steady parallel reference satisfies
\begin{equation}
 (D_{k,w}k_{w,0}')'=E_w-p_0,\qquad
 D_{k,w}(0)k_{w,0}'(0)=0,\qquad k_{w,0}(H)=K.
 \label{mawf:k-bvp}
\end{equation}
Define
\begin{equation}
 Q_0(z)=\int_0^z[E_w(s)-p_0(s)]\,\dd s.
\end{equation}
Two integrations give the unique reference profile
\begin{equation}
 \boxed{
 D_{k,w}k_{w,0}'=Q_0,\qquad
 k_{w,0}(z)=K-\int_z^H\frac{Q_0(s)}{D_{k,w}(s)}\,\dd s.}
 \label{mawf:k-profile}
\end{equation}
The physical upward diffusive flux is \(-Q_0\).
By definition \(Q_0(0)=0\), and the budget gives \(Q_0(H)=0\).
Hence
\begin{equation}
 k_{w,0}(H)=K,\qquad
 k_{w,0}'(0)=k_{w,0}'(H)=0.
 \label{mawf:k-endpoints}
\end{equation}
The ground value is determined by \cref{mawf:k-profile};
an additional \(k_w(0)=0\) generally overdetermines this reference.
The upper zero flux follows from the cutoff budget rather than
being an additional independent boundary condition on a fixed-cutoff
inner problem.

For any sufficiently regular test function \(\phi\),
\begin{equation}
 \boxed{
 -\int_0^H D_{k,w}k_{w,0}'\phi'\,\dd z
 +\int_0^H(p_0-E_w)\phi\,\dd z=0.}
 \label{mawf:weak}
\end{equation}
Indeed, substitution reduces the left side to
\(-[Q_0\phi]_0^H=0\). The constant test function expresses the
budget; for nonconstant modes, ordinary diffusion cancels nonzero
source moments. This is the mechanism required in a spectral/\(hp\)
wall element.

The analytical profile specifies a reference and an initialization;
the evolving \(k_w\) is transported, not reset to this profile.
Any equivalent prescribed redistribution flux from an earlier
flux-form implementation is also removed. The physical flux
\(-D_{k,w}\nabla k_w\) supplies the reference transport.
The cutoff continues to use counterfactual \(p_0(K)\).
For actual production,
\begin{equation}
 \int_0^H(P_A-E_w)\,\dd z=\int_0^H(P_A-p_0)\,\dd z
 \label{mawf:actual-budget}
\end{equation}
need not vanish. Determining the cutoff from actual \(P_A\)
would define a different model.

\subsection{Closed rough-branch formulas and positivity}
\label{mawf:rough}
Assume \(z_v\le0\), and put
\begin{equation}
 d=\nu-az_v\ge\nu,\qquad D_A=d+az,\qquad
 \zeta=\frac{az}{d},\quad L=\frac{aH}{d},\quad
 \beta=\frac{\nu}{d}\in(0,1].
 \label{mawf:rough-variables}
\end{equation}
Then
\begin{align}
 J_0(z)&=\frac1a\ln\frac{d+az}{d},&
 J_2(z)&=\frac1a\left(\frac1d-\frac1{d+az}\right),\\
 p_0(z)&=t_0^2\frac{d+az-\nu}{(d+az)^2},&
 \int_0^zp_0\,\dd s
 &=A_{\eps,m}\left[\ln\frac{d+az}{d}
          -\nu\left(\frac1d-\frac1{d+az}\right)\right].
 \label{mawf:rough-integrals}
\end{align}
Consequently,
\begin{equation}
 \boxed{\begin{gathered}
 \Lambda=\ln(1+L)-\frac{\beta L}{1+L},\qquad
 z_{0,m}=\frac{dH}{aH+d},\\
 z_c=e\,z_{0,m}
 \exp\left[\nu\left(\frac1d-\frac1{aH+d}\right)\right].
 \end{gathered}}
 \label{mawf:rough-cutoff}
\end{equation}
At a prescribed \(z_0<H\), the required rough-branch intercept is
\begin{equation}
 d=\frac{aH z_0}{H-z_0}.
 \label{mawf:d-calibration}
\end{equation}
It must be attainable by the roughness correlation and the selected
wall-layer geometry, normally \(h_r<H\).
Equal viscosity slopes do not remove the finite intercept:
\(D_A(H)=d+aH\), while the outer high-Re viscosity is \(aH\).

Define
\begin{equation}
 R_c(z)=
 \begin{cases}
 z/z_c,&0\le z\le z_c,\\
 1+\ln(z/z_c),&z_c\le z\le H.
 \end{cases}
\end{equation}
The cumulative defect is
\begin{equation}
 \boxed{Q_0(z)=A_{\eps,m}
 \left[R_c(z)-\ln\frac{d+az}{d}
       +\nu\left(\frac1d-\frac1{d+az}\right)\right].}
 \label{mawf:rough-Q}
\end{equation}

\paragraph{Closed kinetic-energy primitive.}
For \(c=az_c/d=L e^{1-\Lambda}\), set
\begin{equation}
 \mathcal F_<(\zeta)=
 \frac{\zeta-\ln(1+\zeta)}c-\frac12\ln^2(1+\zeta)
 +\beta\left[\ln(1+\zeta)+\frac1{1+\zeta}-1\right],
\end{equation}
and
\begin{equation}
 \begin{split}
 \mathcal G(\zeta)
 &=\left[1+\ln\frac{\zeta}{c}+\beta\right]\ln(1+\zeta)
   +\operatorname{Li}_2(-\zeta)
   -\frac12\ln^2(1+\zeta)+\frac{\beta}{1+\zeta},\\
 \operatorname{Li}_2(-\zeta)
 &=-\int_0^\zeta\frac{\ln(1+s)}s\,\dd s\qquad(\zeta\ge0).
 \end{split}
\end{equation}
Join the primitives:
\begin{equation}
 \mathcal F(\zeta)=
 \begin{cases}
 \mathcal F_<(\zeta),&\zeta\le c,\\
 \mathcal F_<(c)+\mathcal G(\zeta)-\mathcal G(c),&\zeta\ge c.
 \end{cases}
\end{equation}
Since
\(\mathcal F'(\zeta)=Q_0/[A_{\eps,m}(1+\zeta)]\),
\begin{equation}
 \boxed{k_{w,0}(z)=K+\frac{A_{\eps,m}}a
  [\mathcal F(az/d)-\mathcal F(aH/d)],\qquad
  \frac{A_{\eps,m}}a=\frac{u_m^2}{\kappa^2}.}
 \label{mawf:k-closed}
\end{equation}
Definite quadrature of \cref{mawf:k-profile} is an equivalent
evaluation. Near \(H\), it can avoid cancellation of close
primitive values.

\paragraph{A sufficient condition for monotone decrease.}
For \(z\ge z_c>0\), direct simplification using
\(A_{\eps,m}=t_0^2/a\) gives
\begin{equation}
 E_w-p_0
 =\frac{A_{\eps,m}[(d+\nu)az+d^2]}{z(d+az)^2}>0.
 \label{mawf:tail-defect}
\end{equation}
Moreover,
\begin{equation}
 p_0'(z)=\frac{a t_0^2[2\nu-(d+az)]}{(d+az)^3}.
 \label{mawf:production-derivative}
\end{equation}

\begin{proposition}[Monotone positive rough reference]
\label{mawf:monotone-proposition}
If \(d\ge2\nu\), \(K>0\), and \(\Lambda\ge1\), then
\begin{equation}
 Q_0(z)<0,\qquad k_{w,0}'(z)<0,\qquad
 k_{w,0}(z)>K>0\quad(0<z<H),
\end{equation}
and \(k_{w,0}(0)>K\).
\end{proposition}
\begin{proof}
The production decreases strictly for \(z>0\).
On the dissipation plateau, \(E_w-p_0\) therefore increases
strictly, and its value at \(z_c\) is positive by
\cref{mawf:tail-defect}. Its zero total integral forces one
negative-to-positive crossing strictly inside \((0,z_c)\).
The cumulative defect is negative until it returns to zero at
\(H\), and \cref{mawf:k-profile} gives the result.
This argument includes \(z_c=H\): the positive endpoint
defect still forces a crossing below \(H\).
\end{proof}
The one-crossing argument works for any positive \(D_{k,w}\).
The following stronger positivity statement uses specifically
\(\sigma_k=1\), hence \(D_{k,w}=D_A\).

\begin{proposition}[Positive reference on the whole rough branch]
\label{mawf:positive-proposition}
Assume the flat measure, constant physical \(\nu>0\),
\(d\ge\nu\), \(K>0\), \(\Lambda\ge1\), \(D_{k,w}=D_A\),
and \(A_{\eps,m}=t_0^2/a\). Then
\begin{equation}
 \boxed{k_{w,0}(z)>K\quad(0\le z<H),\qquad k_{w,0}(H)=K.}
 \label{mawf:positive}
\end{equation}
When \(\nu\le d<2\nu\), the profile can initially increase
before decreasing to \(K\); positivity does not require
monotone decrease.
\end{proposition}
\begin{proof}
Use \(L,\beta\) from \cref{mawf:rough-variables}, and set
\begin{equation}
 \begin{gathered}
 G_\beta(x)=\ln(1+x)-\frac{\beta x}{1+x},\qquad
 c=L e^{1-G_\beta(L)},\\
 \Phi(x;c)=
 \begin{cases}
 x/c-G_\beta(x),&x\le c,\\
 1+\ln(x/c)-G_\beta(x),&x\ge c.
 \end{cases}
 \end{gathered}
\end{equation}
Here \(Q_0=A_{\eps,m}\Phi\), \(\Phi(0)=\Phi(L)=0\), and
\begin{equation}
 k_{w,0}(z)-K=\frac{A_{\eps,m}}a[F(az/d)-F(L)],
 \qquad F(x)=\int_0^x\frac{\Phi(t;c)}{1+t}\,\dd t.
 \label{mawf:positivity-F}
\end{equation}
On the tail,
\(\Phi'=[1+(1+\beta)x]/[x(1+x)^2]>0\), so \(\Phi<0\)
before its endpoint zero. On the plateau,
\[
 \Phi'=\frac1c-p_\beta(x),\qquad
 p_\beta(x)=\frac{1+x-\beta}{(1+x)^2},\qquad
 p_\beta'(x)=\frac{2\beta-1-x}{(1+x)^3}.
\]
Thus \(p_\beta\) decreases or has only one maximum.
The allowed signs of \(\Phi\) are negative throughout the
interior, or positive initially and then negative.
This also holds when \(c=L\), since \(\Phi'(L)>0\).
Consequently \(F\) has no interior minimum.

For fixed \(\beta\), write \(I_\beta(L)=F(L)\) while accounting
for \(c=c(L)\). Since
\[
 \frac{c'(L)}{c(L)}
  =\frac{1+(1+\beta)L}{L(1+L)^2}>0,
\]
and \(\Phi(L)=0\), differentiation gives
\begin{equation}
 I_\beta'(L)=-\frac{c'(L)}{c(L)}
 \left[\frac{c-\ln(1+c)}c+\ln\frac{1+L}{1+c}\right]<0.
\end{equation}
Let \(L_1\) be the unique \(G_\beta(L_1)=1\).
At that lower admissible endpoint \(c=L_1\), and
\(T=\ln(1+L_1)\) satisfies \(1<T<2\). Direct integration gives
\begin{equation}
 I_\beta(L_1)
 =(2-T)\left(\frac{2-T}{2}-\frac T{L_1}\right)<0.
\end{equation}
The sign follows from
\(\ln(1+L_1)>2L_1/(L_1+2)\); the derivative of this difference
is \(L_1^2/[(1+L_1)(L_1+2)^2]>0\).
Hence \(F(L)<0=F(0)\) for every admissible \(L\).
The shape argument now gives \(F(x)>F(L)\) for \(x<L\),
which proves \cref{mawf:positive}.
\end{proof}
The proof covers \(d=\nu,\Lambda=1\) algebraically.
Whether such an endpoint is attainable with the original
roughness correlation and \(h_r<H\) is a separate geometric
admissibility check. The stronger proof is specific to the
flat measure and \(D_{k,w}=D_A\); it is not automatically valid
for another scalar diffusivity or a metric-weighted column.

\paragraph{Why the cutoff gate is not a general positivity test.}
On the smooth branch \(h_r=0\), set \(b=10.7\),
\(x=az/\nu\), \(v=10.7\alpha_m\), and
\(M=aH/\nu\). Let \(y_4=146.409759514\ldots\) solve
\[
 \ln(1+y_4)-\frac{y_4}{1+y_4}=4,
\]
and choose \(M=v+y_4\). Then \(\Lambda=4\) and
\(c=az_c/\nu=M e^{-3}\). With
\[
 B(x)=
 \begin{cases}
 0,&x\le v,\\
 \ln(1+x-v)-(x-v)/(1+x-v),&x>v,
 \end{cases}
 \qquad
 \widehat R_c(x)=
 \begin{cases}x/c,&x\le c,\\1+\ln(x/c),&x\ge c,\end{cases}
\]
the exact ground value is
\begin{equation}
 \frac{k_{w,0}(0)}K
 =1-\frac{1}{\kappa^2q_m}
   \int_0^M\frac{\widehat R_c(x)-B(x)}{1+(x-v)_+}\,\dd x.
 \label{mawf:smooth-counterexample}
\end{equation}
At common \(\kappa=0.41\), this gives
\begin{center}
\begin{tabular}{@{}lrrr@{}}
\toprule
Model&A93&A95&E94/E95\\
\midrule
\(k_{w,0}(0)/K\)&\(-0.371380517\)&\(-0.080424984\)&\(-0.436121275\)\\
\bottomrule
\end{tabular}
\end{center}
The budget is admissible, but the stationary kinetic energy is
inadmissible. Thus outside the rough positivity theorem,
the computed profile needs a separate positivity check.

\subsection{Conservative matching and normal pressure balance}
\label{mawf:matching}
At the flat reference,
\begin{equation}
 U_{w,0}(0)=0,\qquad U_{w,0}(H)=U_{o,0}(H),\qquad
 D_A(H)U_{w,0}'(H)=\nu_{t,o,0}(H)U_{o,0}'(H)=t_0.
 \label{mawf:traction}
\end{equation}
Viscosities and gradients need not be continuous.
The inner Reynolds shear is
\(-t_0\nu_{t,A}/D_A\), while outer Reynolds shear is \(-t_0\);
the inner molecular contribution completes total-traction balance.

Use local transported \(k_w\) in the isotropic inner stress.
With \(\Smat_w=(\nabla\bm U_w+\nabla\bm U_w^T)/2\) and
hydrostatic-reduced pressure \(\pi=p/\rho+gz\), the inner
covariance and total stress are
\begin{equation}
 \R_w=\frac23k_w\Id-2\nu_{t,A}\Smat_w,\qquad
 \bm T_w=-\left(\pi_w+\frac23k_w\right)\Id+2D_A\Smat_w .
 \label{mawf:stress}
\end{equation}
The steady vertical balance and its solution are
\begin{equation}
 \left(\pi_{w,0}+\frac23k_{w,0}\right)'=0,\qquad
 \boxed{\pi_{w,0}(z)=\pi_{o,0}+R_{zz,o,0}-\frac23k_{w,0}(z).}
 \label{mawf:pressure}
\end{equation}
Here \(\pi_{o,0}\) and \(R_{zz,o,0}\) are constant in the
outer reference. Define
\begin{equation}
 F_{93}=\frac{q_{93}^2}{1000+q_{93}^3},\qquad
 B_{95}=
 \frac{q_{95}^2\sqrt{1-9/(2q_{95}^2)}}{1+3q_{95}^2}.
\end{equation}
The original outer stresses give
\begin{equation}
 \boxed{\frac{\pi_{w,0}(H)-\pi_{o,0}}K
 =\begin{cases}
 -17F_{93}/3=-0.0630736219,&\text{A93},\\
 -B_{95}=-0.2710910066,&\text{A95},\\
 0,&\text{E94/E95}.
 \end{cases}}
 \label{mawf:pressure-offset}
\end{equation}
Define \(\Delta\pi_H=\pi_{w,0}(H)-\pi_{o,0}\).
Then \(\pi_{w,0}(z)=\pi_{o,0}+\Delta\pi_H
-\frac23[k_{w,0}(z)-K]\).
Even E94/E95 has a varying inner reduced mechanical pressure,
despite its zero interface offset. A pressure variable that already
includes isotropic turbulent stress must be translated consistently.

The \(k\) value and diffusive flux match:
\begin{equation}
 k_w(H)=k_o(H)=K,\qquad
 D_{k,w}(H)k_{w,0}'(H)
 =\frac{\nu_{t,o,0}(H)}{\sigma_k}k_{o,0}'(H)=0.
 \label{mawf:k-interface}
\end{equation}
For the evolving coupled system these are continuity and conservation
conditions, not independent Dirichlet and Neumann prescriptions on
an isolated outer equation. The outer dissipation equation receives
only
\begin{equation}
 \eps_o(H)=E_w(H)=\frac{A_{\eps,m}}H .
 \label{mawf:eps-interface}
\end{equation}
There is no additional two-sided dissipation-flux condition,
because no inner dissipation equation is solved.
When \(z_c=H\), the inner dissipation derivative is zero and the
outer derivative is \(-A_{\eps,m}/H^2\); their traces still agree.

\subsection{Distinction from the constant-inner-\texorpdfstring{\(k\)}{k} EP-AWF}
\label{mawf:comparison}
At identical sampled \(K\), \(h_r,H,\nu\), compatible coefficients
and cutoff, the constant-inner-\(k\) construction uses
\[
 E_w=\frac{A_{\eps,m}}{\max(z,z_c)},\qquad
 C_k^w=E_w-p_0,\qquad S_{k,w}^{\rm EP}=P_A-p_0.
\]
The modified construction uses
\[
 E_w=\frac{A_{\eps,m}}{\max(z,z_c)},\qquad
 C_k^w=0,\qquad S_{k,w}=P_A-E_w.
\]
Both sources have the same layer integral, but their local
balances differ. At the reference, the first uses prescribed
redistribution to support constant \(k=K\); the second uses
ordinary diffusion of \cref{mawf:k-profile}.
The physical reference transport \(-Q_0\) is the same,
while the evolving models are different.

Setting the host's old switch to \(\chi=0\) does not select this
variant when that switch also restores the original-cutoff
dissipation. The present variant retains the budget cutoff
while removing both the prescribed redistribution source and
its equivalent auxiliary flux. Its scalar policy has no
spatial \(\chi\) fade. The earlier bounded-redistribution
positivity remedy would be a further model change if retained.

\subsection{A specified extension to curved terrain}
\label{mawf:curved}
Use the conservative curved momentum completion with the new
endpoint sample \(K=k|_{\Gamma_H}\), compatible coefficients,
metric-weighted reference budget and \(C_k^w=0\).
The equations are restated here so that the scalar replacement
and roughness-force prescription are explicit.

Let \(\bm n\) point from the ground into the fluid and use the
normal map
\begin{equation}
 \bm x(\bm\xi,r)=\bm X(\bm\xi)+r\bm n(\bm\xi),\qquad
 P=\Id-\bm n\otimes\bm n,\quad B_g=\nabla_\Gamma\bm n,\quad
 A_g=I_T+rB_g,\quad J=\det_TA_g.
 \label{mawf:map}
\end{equation}
The inverses of \(A_g\) act tangentially and are extended by zero
normally. The principal stretches \(1+r\gamma_i\) must be
positive, and the normal map must be globally one-to-one on
\(0\le r\le H(\bm\xi)\). The volume measure is
\(J\,\dd S_\Gamma\,\dd r\). Surface derivatives below are
taken at fixed \(r\) before restriction to the interface.

\paragraph{Reference cutoff in the volume measure.}
Set
\begin{equation}
 E_w(r)=\frac{A_{\eps,m}(K)}{\max(r,r_c)},\qquad
 p_0(r;K)=t_0(K)^2\frac{\nu_{t,A}(r;K)}{D_A(r;K)^2}.
\end{equation}
Use counterfactual zero-forcing production in
\begin{equation}
 \boxed{\int_0^H J(r)E_w(r)\,\dd r
        =\mathcal P_J:=\int_0^H J(r)p_0(r;K)\,\dd r.}
 \label{mawf:metric-budget}
\end{equation}
For \(\mathcal B(c)=A_{\eps,m}\int_0^H J/\max(r,c)\,\dd r\),
\[
 \mathcal B'(c)=-\frac{A_{\eps,m}}{c^2}\int_0^cJ\,\dd r<0,
 \qquad \mathcal B(c)\longrightarrow+\infty\quad(c\downarrow0).
\]
Hence a unique \(0<r_c\le H\) exists exactly when
\begin{equation}
 \boxed{\mathcal P_J\ge
        \frac{A_{\eps,m}}H\int_0^H J\,\dd r.}
 \label{mawf:metric-gate}
\end{equation}
At a fixed footpoint, \(J=(1+\gamma_1r)(1+\gamma_2r)\), and
\begin{equation}
 \frac{\mathcal B(c)}{A_{\eps,m}}
 =1+\ln\frac Hc
 +(\gamma_1+\gamma_2)\left(H-\frac c2\right)
 +\gamma_1\gamma_2\left(\frac{H^2}{2}-\frac{c^2}{6}\right).
 \label{mawf:metric-explicit}
\end{equation}
One monotone scalar root solve replaces the flat exponential.

\paragraph{Actual continuity, momentum, forcing and scalar transport.}
Write \(\bm U_w=\bm u_t+w\bm n\) and
\(\bm v=A_g^{-1}\bm u_t\). With Cartesian velocity components,
\begin{equation}
 \partial_r(Jw)+\operatorname{div}_\Gamma(J\bm v)=0,
 \label{mawf:continuity}
\end{equation}
\begin{equation}
 \begin{split}
 J\partial_t\bm U_w
 &+\partial_r\{J(w\bm U_w-T_w\bm n)\}\\
 &+\operatorname{Div}_\Gamma
   \{J(\bm U_w\otimes\bm v-T_wA_g^{-T})\}=J\bm f_w .
 \end{split}
 \label{mawf:curved-momentum}
\end{equation}
Here \(T_w\) is \cref{mawf:stress};
the outer stress is \(T_o=-\pi_o\Id-\R_m\).
The operator \(\operatorname{Div}_\Gamma\) takes the surface
divergence of each Cartesian row. These fluxes already contain
the geometric effects of stress and pressure.

Retain the explicit roughness-force completion
\begin{equation}
 \begin{split}
 a_r(r)&=\mathbf1_{\{r>h_r\}},\qquad
 \bm f_w=\bm f_\pi+(1-a_r)\bm F_A,\\
 \bm F_A&=\mathcal E_H^+P
 \left[\operatorname{Div}_r(\bm u_{o,t}\otimes\bm u_{o,t})
             +\nabla_r\pi_o-\bm f_\pi\right],
 \qquad \bm u_{o,t}=P\bm U_o .
 \end{split}
 \label{mawf:rough-force}
\end{equation}
Here \(\bm f_\pi\) is the remaining body acceleration after the
same hydrostatic-pressure convention used outside,
\(\nabla_rf=A_g^{-T}\nabla_\Gamma f\), and
\[
 \operatorname{Div}_r Q
 =J^{-1}\operatorname{Div}_\Gamma(JQA_g^{-T})
 \quad\text{for }Q\bm n=0.
\]
The sampler \(\mathcal E_H^+\) takes one-sided outer values after
derivatives at fixed \(r\) have been evaluated; extend the resulting
\(\bm F_A\) constantly along each normal ray.
Its work has no guaranteed drag sign and must be retained in
energy accounting.

The flat scalar reconstruction is the declared quasi-steady
parallel-flow reduction of this completion: inner storage,
normal advection, lateral viscous transport and the relevant
\(\nabla_\parallel w\) contributions are neglected, and the
remaining convection/pressure/body-force balance is replaced
by the height-constant sample \(\bm F_A\).
This also replaces the inner effective-pressure gradient by the
sampled outer mechanical-pressure gradient; the discarded
\(\nabla_\parallel(\pi_w+2k_w/3-\pi_o)\) is generally nonzero.
Under those approximations
\((D_A\bm U_\parallel')'=a_r\bm F_A\).
For the curved or unsteady patch, use the full equation
\cref{mawf:curved-momentum}, without additionally imposing that
scalar reduction. All discarded terms vanish on the homogeneous
zero-forcing reference.

The modified actual inner \(k\) equation is
\begin{equation}
 \boxed{\begin{aligned}
 J\partial_t k_w
 &+\partial_r\{J(wk_w-D_{k,w}\partial_rk_w)\}\\
 &+\operatorname{div}_\Gamma
 \{J[\bm v k_w-D_{k,w}A_g^{-1}A_g^{-T}\nabla_\Gamma k_w]\}
 =J(P_A-E_w).
 \end{aligned}}
 \label{mawf:curved-k}
\end{equation}
For incompressible flow,
\[
 P_A=2\nu_{t,A}\Smat_w:\Smat_w .
\]
The isotropic stress uses local \(k_w\), while the viscosity
uses sampled \(K\). If incompressibility is enforced only weakly,
the production contraction \(-\R_w:\nabla\bm U_w\) also contains
\(-(2k_w/3)\nabla\cdot\bm U_w\); its treatment must agree with
the discrete stress-work and pressure convention.
The base-normal tangential traction is
\begin{equation}
 \bm t_A=D_A\left[
 \partial_r\bm u_t+A_g^{-T}\nabla_\Gamma w-B_gA_g^{-1}\bm u_t
 \right].
 \label{mawf:curved-traction}
\end{equation}
Curvature, normal motion and normal pressure balance are thus
retained in the full patch.

The actual source budget remains
\[
 \int_0^HJ(P_A-E_w)\,\dd r
 =\int_0^HJ(P_A-p_0)\,\dd r,
\]
which need not vanish. A frozen reference column can be defined by
\begin{equation}
 Q_J(r)=\int_0^rJ(s)[E_w(s)-p_0(s)]\,\dd s,\qquad
 k_{\rm col}(r)=K-\int_r^H
        \frac{Q_J(s)}{J(s)D_{k,w}(s)}\,\dd s.
 \label{mawf:curved-column}
\end{equation}
Its base-normal endpoint fluxes vanish.
The one-crossing proof for \(d\ge2\nu\) extends to this column
for positive \(J,D_{k,w}\).
The stronger whole-rough-branch proof is not automatically
metric weighted. The column balances only the reference source
and normal diffusion; it is not generally an equilibrium of
the complete terrain-flow equations.

\paragraph{Actual interface geometry and shared fluxes.}
For variable \(H\), put
\begin{equation}
 \bm q_H=A_{g,H}^{-T}\nabla_\Gamma H,\qquad
 \beta_H=\sqrt{1+|\bm q_H|^2},\qquad
 \bm N_H=\frac{\bm n-\bm q_H}{\beta_H}.
 \label{mawf:interface-normal}
\end{equation}
Use common velocity and total traction,
\(\bm U_w=\bm U_o=\widehat{\bm U}\) and
\(T_w\bm N_H=T_o\bm N_H\), with one conservative momentum
flux and opposite orientations on the two sides.
The mass flux per unit ground area is
\(m_H=J_H(w-\bm u_t\cdot\bm q_H)\).
At a common \(k\) trace \(\widehat k\), match
\begin{equation}
 J_H\beta_H(\bm U k-D_k\nabla k)\cdot\bm N_H
 =m_H\widehat k-J_HD_k
 \left[\partial_rk-\bm q_H\cdot A_{g,H}^{-T}\nabla_\Gamma k\right].
 \label{mawf:curved-k-flux}
\end{equation}
Zero base-normal column flux alone does not establish zero flux
through a tilted interface. Surface differentiation of the trace
obeys
\(\nabla_\Gamma K=(\nabla_\Gamma k_w)|_H+
(\partial_rk_w)|_H\nabla_\Gamma H\).
The outer dissipation trace is \(E_w(H)=A_{\eps,m}/H\).
Full normal-traction matching replaces the flat pressure offset;
there is still no independent inner dissipation-flux prescription.

\subsection{Reference realizability, positivity, and stability limits}
\label{mawf:stability}
On a positive rough reference the smallest inner covariance
eigenvalue is
\begin{equation}
 \lambda_{\min}(\R_{w,0})
 =\frac23k_{w,0}-t_0\frac{\nu_{t,A}}{D_A}
 >K\left(\frac23-\frac1{q_m}\right)>0.
 \label{mawf:covariance}
\end{equation}
All three \(q_m\) exceed \(3/2\). Their original outer reference
covariances are also positive definite. These are reference-state
results, not global guarantees for forced or separated states.

With endpoint \(K\), coefficients, cutoff and reference production
held fixed, a wall-normal perturbation \(v=k_w-k_{w,0}\) obeys
\[
 v_t=(D_{k,w}v')',\qquad v(H)=0,\qquad D_{k,w}v'(0)=0,
\]
and
\begin{equation}
 \frac12\frac{\dd}{\dd t}\int_0^H v^2\,\dd z
 =-\int_0^H D_{k,w}(v')^2\,\dd z\le0.
 \label{mawf:scalar-stability}
\end{equation}
This conditional scalar estimate is separate from full coupled
stability.

Dynamic nonnegativity of actual \(k_w\) is not automatic:
at a flat interior zero minimum, positive endpoint \(K\), and
locally zero production and diffusion,
\[
 \partial_t k_w=-E_w(K)<0.
\]
Thus the positive rough reference does not make the evolving
sampled-\(K\) sink positivity preserving. An additional source
modification or numerical positivity device must be identified
and checked separately if introduced.

The wall changes do not alter the original outer constitutive
domains or incremental momentum response. A93 can lose covariance
realizability away from the reference, A95 retains its real-radicand
restriction, and positive covariance in E94/E95 does not ensure
positive incremental momentum diffusion. The previously established
outer constitutive restrictions therefore remain relevant.

\subsection{Illustrative references and implementation implications}
\label{mawf:verification}
Take
\begin{equation}
 \ust=0.4\ \mathrm{m\,s^{-1}},\quad
 \nu=1.5\times10^{-5}\ \mathrm{m^2\,s^{-1}},\quad
 h_r=0.1\ \mathrm m,\quad H=2\ \mathrm m,\quad z_K=z_m=H.
 \label{mawf:inputs}
\end{equation}
Using the original roughness correlation gives
\cref{mawf:numerical-table}. The same sand-grain input is retained,
so aerodynamic roughness is an output rather than a common
prescribed input.

\begin{table}[htbp]
\centering\small
\caption{Analytical rough references. \(K,k(0)\) are in
\(\mathrm{m^2\,s^{-2}}\), and lengths are in metres.}
\label{mawf:numerical-table}
\begin{tabular}{@{}llrrrrr@{}}
\toprule
Model&\(\kappa\) choice&\(K\)&\(k(0)\)&\(\Lambda\)&\(z_c\)&\(z_{0,m}\)\\
\midrule
A93&original&0.54420247&1.08630845&5.90564847&0.01480928&0.00536024\\
A95&original&0.61152007&1.22075495&5.89870963&0.01491240&0.00539286\\
E94/E95&original&0.53192103&1.11770080&5.90535573&0.01481362&0.00535834\\
\midrule
A93&common \(0.41\)&0.54420247&1.13595953&5.90420900&0.01483061&0.00536413\\
A95&common \(0.41\)&0.61152007&1.20319430&5.89920134&0.01490507&0.00539152\\
E94/E95&common \(0.41\)&0.53192103&1.12369422&5.90518519&0.01481614&0.00535880\\
\bottomrule
\end{tabular}
\end{table}

All six cases satisfy \(d\ge2\nu\) and \(\Lambda>1\), so their
positive monotone profiles have an analytical guarantee.
The interface pressure offsets are
\(-0.034324820606\), \(-0.165777590815\), and \(0\)
\(\mathrm{m^2\,s^{-2}}\), respectively.
At fixed \(\ust\) these offsets do not change with \(\kappa\),
because \(K=q_m\ust^2\) is unchanged; the inner pressure profile
does change with \(k_w(z)\).

Independent quadrature of the production and dissipation budgets,
the cumulative defect, and the closed primitive agrees with the
analytical expressions. Separately integrated source and diffusion
moments for Legendre degrees \(0,\ldots,6\) give a maximum absolute
weak residual of \(2.33\times10^{-16}\) across these cases;
closed \(k\) versus quadrature differs by at most
\(6.89\times10^{-15}\) in SI units. These are checks of the
formulas, not CFD or terrain validation.

The model-specific \(\kappa\) affects the viscosity slope,
dissipation amplitude, cutoff and inner profile. A further useful
identity is
\[
 \frac{\kappa^2[k_{w,0}(z)-K]}{u_m^2}
 =\mathcal F(az/d)-\mathcal F(aH/d).
\]
At common \(\kappa,\ust,H,z_0,\nu\), if each roughness input
is calibrated to the same \(z_0\), \cref{mawf:d-calibration}
makes \(a,d\) identical. Then the zero-forcing wall velocity,
production, sink and cumulative defect coincide across the models,
and their \(k\) profiles differ only by \(K=q_m\ust^2\).
Their forced responses can still differ through the calibrated
\(h_r\) in the forcing cutoff.

For a spectral/\(hp\) implementation, the exact continuum weak
identity \cref{mawf:weak} does not by itself make a finite polynomial
interpolant stationary. The ordinary diffusive flux, nonzero source
moments and shared interface terms must be represented and integrated
consistently. A local discrete wall solve using the same source,
quadrature and boundary contract can supply a stationary discrete
profile. Checking only the cell-mean energy budget is insufficient.

The selected variant therefore consists of endpoint sampling,
the two compatible coefficients, the budget cutoff and its plateau,
retained finite-\(\nu\) transport with \(P_A-E_w\), and conservative
velocity/stress/scalar coupling. Its exactness claim is the
flat neutral reference of these specified equations. Positivity
outside the rough theorem, global coupled stability and fully
discrete preservation remain separate requirements.

% The main note already calls \appendix; do not call it again here.
% Uses the main note's amsmath/amssymb/amsthm/bm/booktabs/natbib/cleveref
% packages, proposition environment, and \eps, \ust, \dd, \Id, \R, \Smat.
% Bibliography keys: Shih1993, Shih1995ASM, Shih1994EVM, Shih1995EVM,
% and SCI2006. All labels introduced here have the unique prefix mawf:.
% No document class, document environment, or bibliography is added.

\section{AWF with coefficients compatible with A93, A95 and E94/E95,
a budget-determined cutoff, and a nonconstant inner kinetic-energy profile}
\label{mawf:appendix}

This appendix specifies a modified analytical wall function for each of
the three outer closures A93, A95 and E94/E95
\citep{Shih1993,Shih1995ASM,Shih1994EVM,Shih1995EVM}.
The objective is to retain its exact neutral outer equilibrium while
allowing kinetic energy to vary inside the modeled wall layer.
Two auxiliary coefficients are matched to the selected outer branch;
the dissipation cutoff is determined by the reference energy budget;
and the prescribed redistribution source is removed. Ordinary diffusion
of the transported inner kinetic energy then balances the local
production--dissipation difference.

The rough-wall viscosity kernel and analytical momentum reconstruction
start from the AWF of Suga, Craft and Iacovides \citep{SCI2006}.
The compatible coefficients, endpoint sampling, budget cutoff and
variable inner reference specified below are explicit modifications.
Selecting this variant replaces the earlier wall-coefficient,
sampling and \(\chi\)-dependent scalar-source choices. Its own
dissipation law is retained with \(C_k^w=0\).

\subsection{Scope, endpoint sampling, and compatible coefficients}
\label{mawf:coefficients}

Use the selected high-Re outer equations: molecular momentum and scalar
diffusion are omitted from the outer analytical reference, and the
E94/E95 dissipation denominator is reduced from
\(k+\sqrt{\nu\eps}\) to \(k\) in that outer limit.
The inner AWF retains the fixed physical \(\nu>0\) in its momentum,
production and scalar-diffusion formulas. These are two declared
equation regimes; the result is not an exact logarithmic solution of
an unchanged finite-Re outer system.

Only inner \(k_w\) is transported. Inner dissipation is a dependent
total sink, and no inner dissipation transport equation is introduced.
The inner stress uses the auxiliary linear viscosity specified below;
the original A93/A95 nonlinear stress laws remain the outer closures.
Coordinates are \(x\) streamwise, \(y\) spanwise, and \(z\) vertical.
All stresses are per unit density.

Let \(m\in\{93,95,E\}\) identify the outer model.
Its neutral equilibrium ratio \(q_m=k_{o,0}/\ust^2\) is
\begin{equation}
 \boxed{
 q_{93}=\frac{57+\sqrt{6249}}{40},\qquad
 q_{95}=\frac{3+\sqrt{61}}{2\sqrt2},\qquad
 q_E=\frac{3+\sqrt{41}}{2\sqrt2}.}
 \label{mawf:q}
\end{equation}
These roots follow from \(C_{\mu,m}(q_m)q_m^2=1\) in simple shear.
The nonlinear corrections of A93 and A95 have zero \(xz\) component
there, but the full Reynolds shear is \(R_{xz,o,0}=-\ust^2\).

The outer dissipation equation additionally requires
\begin{equation}
 \boxed{\kappa^2q_m=\sigma_{\eps,m}\Delta_m,\qquad
 \Delta_{93}=\Delta_{95}=0.48,\qquad
 \Delta_E=1.9-0.43q_E.}
 \label{mawf:kappa}
\end{equation}
For E94/E95 the branch \(C_1=0.43\) is active at \(q_E\).
Retaining the original outer coefficients means selecting
\(\kappa_m=\sqrt{\sigma_{\eps,m}^{\rm orig}\Delta_m/q_m}\).
A prescribed common \(\kappa\) instead requires the single outer
replacement \(\sigma_{\eps,m}^{\rm eq}=\kappa^2q_m/\Delta_m\).
The exact values implied by the printed coefficients are listed in
\cref{mawf:outer-table}; they need not coincide with a nominal
\(\kappa\) used in the original empirical calibration.

\begin{table}[htbp]
\centering\small
\caption{Outer choices. The last column is the required one-coefficient
change when a common \(\kappa=0.41\) is selected.}
\label{mawf:outer-table}
\begin{tabular}{@{}lrrrr@{}}
\toprule
Model & \(q_m\) & original \(\sigma_\eps\)& compatible \(\kappa_m\)&
\(\sigma_\eps^{\,0.41}\)\\
\midrule
A93 &3.401265417&1.3&0.428323629&1.191151493\\
A95 &3.822000426&1.3&0.404060995&1.338496399\\
E94/E95 &3.324506456&1.2&0.412087472&1.187873345\\
\bottomrule
\end{tabular}
\end{table}

Use endpoint sampling and tail matching:
\begin{equation}
 \boxed{z_K=z_m=H,\qquad K=k_w(H)=k_o(H)>0.}
 \label{mawf:sampling}
\end{equation}
The subscripts \(w,o\) denote wall layer and outer region; subscript
\(0\) denotes their steady flat reference. At each positive current
sample \(K\), define counterfactual zero-forcing scales
\begin{equation}
 u_m(K)=\sqrt{\frac K{q_m}},\qquad
 t_0(K)=\frac K{q_m},\qquad
 A_{\eps,m}(K)=\frac{u_m(K)^3}{\kappa}.
 \label{mawf:scales}
\end{equation}
At the chosen inlet reference \(K=q_m\ust^2\), \(u_m=\ust\).
Away from that reference, \(t_0\) is not imposed as the actual
wall traction. Model and sample arguments are suppressed below.

Match the wall dissipation-tail amplitude and viscosity slope:
\begin{equation}
 \frac{K^{3/2}}{c_{\ell,m}^w}=\frac{u_m^3}{\kappa},\qquad
 c_{\mu,m}^wc_{\ell,m}^w\sqrt K=\kappa u_m.
\end{equation}
The two compatible auxiliary coefficients and their product are
\begin{equation}
 \boxed{c_{\ell,m}^w=\kappa q_m^{3/2},\qquad
 c_{\mu,m}^w=q_m^{-2},\qquad
 \alpha_m=c_{\mu,m}^wc_{\ell,m}^w=\frac{\kappa}{\sqrt{q_m}}.}
 \label{mawf:wall-coefficients}
\end{equation}
Consequently, with \(a=\alpha_m\sqrt K\),
\begin{equation}
 \boxed{a=\kappa u_m,\qquad
 \frac{K^{3/2}}{c_{\ell,m}^w}=A_{\eps,m}=\frac{t_0^2}{a}.}
 \label{mawf:identities}
\end{equation}
These two matching requirements uniquely determine the two auxiliary
coefficients. Traction continuity by itself does not require
equal inner and outer viscosity slopes.
The auxiliary \(c_\mu^w\) is constant for the selected branch;
it is not the current gradient-dependent outer \(C_\mu\).

\begin{table}[htbp]
\centering\small
\caption{Compatible wall coefficients. Each row uses the same
\(\kappa\) in the outer equations and the wall model.}
\label{mawf:coefficient-table}
\begin{tabular}{@{}llrrrr@{}}
\toprule
Model&outer choice&\(\kappa\)&\(c_\mu^w\)&\(c_\ell^w\)&\(\alpha\)\\
\midrule
A93&original&0.428323629&0.086440835&2.686784401&0.232247887\\
A95&original&0.404060995&0.068457106&3.019138428&0.206681480\\
E94/E95&original&0.412087472&0.090478551&2.497929531&0.226009045\\
\midrule
A93&common \(\kappa\)&0.410000000&0.086440835&2.571844117&0.222312353\\
A95&common \(\kappa\)&0.410000000&0.068457106&3.063514593&0.209719344\\
E94/E95&common \(\kappa\)&0.410000000&0.090478551&2.485276011&0.224864173\\
\bottomrule
\end{tabular}
\end{table}

Retain the original roughness-intercept correlation. With sand-grain
roughness \(h_r\ge0\), define, for \(h_r>0\),
\begin{equation}
 h^*=\frac{h_r\sqrt K}{\nu},\quad x=\frac{h^*}{70},\quad
 m_r=\max(0.5-0.4x^{0.7},\,1-0.79x^{-0.28}),\quad
 b=10.7(1-x^{m_r});
 \label{mawf:roughness}
\end{equation}
for \(h_r=0\), set \(b=10.7\). The wall coefficients are
\begin{equation}
 \begin{gathered}
 z_v=\frac{\nu b}{\sqrt K},\qquad
 \nu_{t,A}(z)=a(z-z_v)_+,\qquad D_A(z)=\nu+\nu_{t,A}(z),\\
 D_{k,w}(z)=\nu+\frac{\nu_{t,A}(z)}{\sigma_k}=D_A(z).
 \end{gathered}
 \label{mawf:viscosity}
\end{equation}
Here \((r)_+=\max(r,0)\), and the last equality uses \(\sigma_k=1\)
for all three models. These coefficients sample endpoint \(K\);
local transported \(k_w(z)\) is allowed to vary. Substituting local
\(k_w\) into the kernel would define a different model. Likewise,
an interior sample or average cannot simply replace the endpoint
sample once the inner reference is nonconstant.

\subsection{Momentum, production, and the outer velocity anchor}
\label{mawf:momentum}
For a flat quasi-steady column, put
\begin{equation}
 L_h(z)=(z-h_r)_+,\quad
 J_0(z)=\int_0^z\frac{\dd s}{D_A(s)},\quad
 J_1(z)=\int_0^z\frac{L_h(s)}{D_A(s)}\,\dd s,\quad
 J_2(z)=\int_0^z\frac{\dd s}{D_A(s)^2}.
 \label{mawf:J}
\end{equation}
With prescribed constant tangential forcing \(F\), retain
\begin{equation}
 (D_AU_A')'=F\,\mathbf1_{\{z>h_r\}},\qquad
 U_A(0)=0,\quad U_A(H)=U_H.
 \label{mawf:momentum-bvp}
\end{equation}
Integration yields the original analytical form, now evaluated
with the compatible viscosity:
\begin{equation}
 \boxed{\begin{aligned}
 t_A&=\frac{U_H-FJ_1(H)}{J_0(H)},&
 U_A(z)&=t_AJ_0(z)+FJ_1(z),\\
 P_A(z)&=\nu_{t,A}(z)
       \left|\frac{t_A+FL_h(z)}{D_A(z)}\right|^2 .
 \end{aligned}}
 \label{mawf:forced}
\end{equation}
The upper traction is \(T_H=t_A+FL_h(H)\), generally different
from ground traction \(t_A\). The vector tangential form uses
componentwise reconstruction and the Euclidean square in \(P_A\).
The coordinates remain the physical Cartesian coordinates.

At the counterfactual zero-forcing reference choose
\begin{equation}
 U_{H,0}=t_0J_0(H),\qquad U_{w,0}(z)=t_0J_0(z),\qquad
 p_0(z)=t_0^2\frac{\nu_{t,A}(z)}{D_A(z)^2}.
 \label{mawf:reference-production}
\end{equation}
Its integrated production is exactly
\begin{equation}
 \boxed{\mathcal P_0=\int_0^H p_0\,\dd z
     =t_0^2[J_0(H)-\nu J_2(H)].}
 \label{mawf:production-budget}
\end{equation}
Here \(p_0\) is pointwise reference production and
\(\mathcal P_0/H\) is its height mean. Neither is the actual
production of a general current flow. The molecular part of mean
viscous dissipation is excluded from turbulent production; hence
the term \(-\nu J_2\) must be retained.

For \(z\ge H\), the compatible outer reference is
\begin{equation}
 \boxed{\begin{aligned}
 U_{o,0}(z)&=t_0J_0(H)+\frac{u_m}{\kappa}\ln\frac zH,\\
 k_{o,0}(z)&=K,\qquad
 \eps_{o,0}(z)=\frac{A_{\eps,m}}z,\qquad
 \nu_{t,o,0}(z)=az .
 \end{aligned}}
 \label{mawf:outer-reference}
\end{equation}
Its aerodynamic intercept is
\begin{equation}
 z_{0,m}=H\exp[-\kappa u_mJ_0(H)]=H\exp[-aJ_0(H)].
 \label{mawf:z0}
\end{equation}
Keeping the same \(h_r\) after changing coefficients generally
changes \(z_{0,m}\). A separately prescribed \(z_0\) requires
\begin{equation}
 J_0(H)=\frac{\ln(H/z_0)}{\kappa\ust}
 \label{mawf:z0-calibration}
\end{equation}
at the selected inlet reference. A physical roughness input calibrated
there is subsequently held fixed as the endpoint sample evolves.

\subsection{Budget-determined dissipation and admissibility}
\label{mawf:cutoff-section}
Use a positive continuous dependent sink and no prescribed redistribution:
\begin{equation}
 \boxed{E_w(z)=\frac{A_{\eps,m}}{\max(z,z_c)},\qquad
 C_k^w(z)=0,\qquad 0<z_c\le H.}
 \label{mawf:dissipation}
\end{equation}
Its integral and upper trace are
\begin{equation}
 \int_0^HE_w\,\dd z=A_{\eps,m}\left(1+\ln\frac H{z_c}\right),
 \qquad E_w(H)=\frac{A_{\eps,m}}H.
 \label{mawf:dissipation-integral}
\end{equation}
Zero ground \(k\) flux and the constant outer reference require
\(\int_0^H(E_w-p_0)\,\dd z=0\). Thus
\begin{equation}
 \boxed{\Lambda=\frac{\mathcal P_0}{A_{\eps,m}}
       =a[J_0(H)-\nu J_2(H)],\qquad
 z_c=H\exp(1-\Lambda).}
 \label{mawf:cutoff}
\end{equation}
The integral in \cref{mawf:dissipation-integral} decreases strictly
with \(z_c\), tends to infinity as \(z_c\downarrow0\), and equals
\(A_{\eps,m}\) at \(z_c=H\). Therefore
\begin{equation}
 \boxed{\text{a unique }0<z_c\le H\text{ exists}
        \quad\Longleftrightarrow\quad \Lambda\ge1.}
 \label{mawf:gate}
\end{equation}
If the gate fails, this sink family cannot satisfy both the budget
and the outer trace. Clipping \(z_c\) to \(H\) loses the budget.

The original molecular cutoff
\(z_{\eps,A}=2c_\ell^w\nu/\sqrt K\) is replaced by
\cref{mawf:cutoff}. The plateau is \(A_{\eps,m}/z_c\).
Imposing also \(E_w(0)=2\nu K/z_c^2\) would force the original
molecular cutoff again. The replacement cutoff and plateau are
explicit modifications of the dissipation law.
The retained finite-viscosity contribution also gives
\begin{equation}
 \boxed{z_c=e\,z_{0,m}\exp[a\nu J_2(H)]>e\,z_{0,m},}
 \label{mawf:cutoff-z0}
\end{equation}
and the mechanical-work budget
\begin{equation}
 t_0U_{H,0}
 =\nu t_0^2J_2(H)+\int_0^Hp_0\,\dd z
 =\nu t_0^2J_2(H)+\int_0^HE_w\,\dd z.
 \label{mawf:work}
\end{equation}

\subsection{Nonconstant inner kinetic energy and every weak test mode}
\label{mawf:inner-k}
Retain the actual inner transport equation
\begin{equation}
 \partial_tk_w+\bm U_w\cdot\nabla k_w
 =\nabla\cdot(D_{k,w}\nabla k_w)+P_A-E_w.
 \label{mawf:k-transport}
\end{equation}
The steady parallel reference satisfies
\begin{equation}
 (D_{k,w}k_{w,0}')'=E_w-p_0,\qquad
 D_{k,w}(0)k_{w,0}'(0)=0,\qquad k_{w,0}(H)=K.
 \label{mawf:k-bvp}
\end{equation}
Define
\begin{equation}
 Q_0(z)=\int_0^z[E_w(s)-p_0(s)]\,\dd s.
\end{equation}
Two integrations give the unique reference profile
\begin{equation}
 \boxed{
 D_{k,w}k_{w,0}'=Q_0,\qquad
 k_{w,0}(z)=K-\int_z^H\frac{Q_0(s)}{D_{k,w}(s)}\,\dd s.}
 \label{mawf:k-profile}
\end{equation}
The physical upward diffusive flux is \(-Q_0\).
By definition \(Q_0(0)=0\), and the budget gives \(Q_0(H)=0\).
Hence
\begin{equation}
 k_{w,0}(H)=K,\qquad
 k_{w,0}'(0)=k_{w,0}'(H)=0.
 \label{mawf:k-endpoints}
\end{equation}
The ground value is determined by \cref{mawf:k-profile};
an additional \(k_w(0)=0\) generally overdetermines this reference.
The upper zero flux follows from the cutoff budget rather than
being an additional independent boundary condition on a fixed-cutoff
inner problem.

For any sufficiently regular test function \(\phi\),
\begin{equation}
 \boxed{
 -\int_0^H D_{k,w}k_{w,0}'\phi'\,\dd z
 +\int_0^H(p_0-E_w)\phi\,\dd z=0.}
 \label{mawf:weak}
\end{equation}
Indeed, substitution reduces the left side to
\(-[Q_0\phi]_0^H=0\). The constant test function expresses the
budget; for nonconstant modes, ordinary diffusion cancels nonzero
source moments. This is the mechanism required in a spectral/\(hp\)
wall element.

The analytical profile specifies a reference and an initialization;
the evolving \(k_w\) is transported, not reset to this profile.
Any equivalent prescribed redistribution flux from an earlier
flux-form implementation is also removed. The physical flux
\(-D_{k,w}\nabla k_w\) supplies the reference transport.
The cutoff continues to use counterfactual \(p_0(K)\).
For actual production,
\begin{equation}
 \int_0^H(P_A-E_w)\,\dd z=\int_0^H(P_A-p_0)\,\dd z
 \label{mawf:actual-budget}
\end{equation}
need not vanish. Determining the cutoff from actual \(P_A\)
would define a different model.

\subsection{Closed rough-branch formulas and positivity}
\label{mawf:rough}
Assume \(z_v\le0\), and put
\begin{equation}
 d=\nu-az_v\ge\nu,\qquad D_A=d+az,\qquad
 \zeta=\frac{az}{d},\quad L=\frac{aH}{d},\quad
 \beta=\frac{\nu}{d}\in(0,1].
 \label{mawf:rough-variables}
\end{equation}
Then
\begin{align}
 J_0(z)&=\frac1a\ln\frac{d+az}{d},&
 J_2(z)&=\frac1a\left(\frac1d-\frac1{d+az}\right),\\
 p_0(z)&=t_0^2\frac{d+az-\nu}{(d+az)^2},&
 \int_0^zp_0\,\dd s
 &=A_{\eps,m}\left[\ln\frac{d+az}{d}
          -\nu\left(\frac1d-\frac1{d+az}\right)\right].
 \label{mawf:rough-integrals}
\end{align}
Consequently,
\begin{equation}
 \boxed{\begin{gathered}
 \Lambda=\ln(1+L)-\frac{\beta L}{1+L},\qquad
 z_{0,m}=\frac{dH}{aH+d},\\
 z_c=e\,z_{0,m}
 \exp\left[\nu\left(\frac1d-\frac1{aH+d}\right)\right].
 \end{gathered}}
 \label{mawf:rough-cutoff}
\end{equation}
At a prescribed \(z_0<H\), the required rough-branch intercept is
\begin{equation}
 d=\frac{aH z_0}{H-z_0}.
 \label{mawf:d-calibration}
\end{equation}
It must be attainable by the roughness correlation and the selected
wall-layer geometry, normally \(h_r<H\).
Equal viscosity slopes do not remove the finite intercept:
\(D_A(H)=d+aH\), while the outer high-Re viscosity is \(aH\).

Define
\begin{equation}
 R_c(z)=
 \begin{cases}
 z/z_c,&0\le z\le z_c,\\
 1+\ln(z/z_c),&z_c\le z\le H.
 \end{cases}
\end{equation}
The cumulative defect is
\begin{equation}
 \boxed{Q_0(z)=A_{\eps,m}
 \left[R_c(z)-\ln\frac{d+az}{d}
       +\nu\left(\frac1d-\frac1{d+az}\right)\right].}
 \label{mawf:rough-Q}
\end{equation}

\paragraph{Closed kinetic-energy primitive.}
For \(c=az_c/d=L e^{1-\Lambda}\), set
\begin{equation}
 \mathcal F_<(\zeta)=
 \frac{\zeta-\ln(1+\zeta)}c-\frac12\ln^2(1+\zeta)
 +\beta\left[\ln(1+\zeta)+\frac1{1+\zeta}-1\right],
\end{equation}
and
\begin{equation}
 \begin{split}
 \mathcal G(\zeta)
 &=\left[1+\ln\frac{\zeta}{c}+\beta\right]\ln(1+\zeta)
   +\operatorname{Li}_2(-\zeta)
   -\frac12\ln^2(1+\zeta)+\frac{\beta}{1+\zeta},\\
 \operatorname{Li}_2(-\zeta)
 &=-\int_0^\zeta\frac{\ln(1+s)}s\,\dd s\qquad(\zeta\ge0).
 \end{split}
\end{equation}
Join the primitives:
\begin{equation}
 \mathcal F(\zeta)=
 \begin{cases}
 \mathcal F_<(\zeta),&\zeta\le c,\\
 \mathcal F_<(c)+\mathcal G(\zeta)-\mathcal G(c),&\zeta\ge c.
 \end{cases}
\end{equation}
Since
\(\mathcal F'(\zeta)=Q_0/[A_{\eps,m}(1+\zeta)]\),
\begin{equation}
 \boxed{k_{w,0}(z)=K+\frac{A_{\eps,m}}a
  [\mathcal F(az/d)-\mathcal F(aH/d)],\qquad
  \frac{A_{\eps,m}}a=\frac{u_m^2}{\kappa^2}.}
 \label{mawf:k-closed}
\end{equation}
Definite quadrature of \cref{mawf:k-profile} is an equivalent
evaluation. Near \(H\), it can avoid cancellation of close
primitive values.

\paragraph{A sufficient condition for monotone decrease.}
For \(z\ge z_c>0\), direct simplification using
\(A_{\eps,m}=t_0^2/a\) gives
\begin{equation}
 E_w-p_0
 =\frac{A_{\eps,m}[(d+\nu)az+d^2]}{z(d+az)^2}>0.
 \label{mawf:tail-defect}
\end{equation}
Moreover,
\begin{equation}
 p_0'(z)=\frac{a t_0^2[2\nu-(d+az)]}{(d+az)^3}.
 \label{mawf:production-derivative}
\end{equation}

\begin{proposition}[Monotone positive rough reference]
\label{mawf:monotone-proposition}
If \(d\ge2\nu\), \(K>0\), and \(\Lambda\ge1\), then
\begin{equation}
 Q_0(z)<0,\qquad k_{w,0}'(z)<0,\qquad
 k_{w,0}(z)>K>0\quad(0<z<H),
\end{equation}
and \(k_{w,0}(0)>K\).
\end{proposition}
\begin{proof}
The production decreases strictly for \(z>0\).
On the dissipation plateau, \(E_w-p_0\) therefore increases
strictly, and its value at \(z_c\) is positive by
\cref{mawf:tail-defect}. Its zero total integral forces one
negative-to-positive crossing strictly inside \((0,z_c)\).
The cumulative defect is negative until it returns to zero at
\(H\), and \cref{mawf:k-profile} gives the result.
This argument includes \(z_c=H\): the positive endpoint
defect still forces a crossing below \(H\).
\end{proof}
The one-crossing argument works for any positive \(D_{k,w}\).
The following stronger positivity statement uses specifically
\(\sigma_k=1\), hence \(D_{k,w}=D_A\).

\begin{proposition}[Positive reference on the whole rough branch]
\label{mawf:positive-proposition}
Assume the flat measure, constant physical \(\nu>0\),
\(d\ge\nu\), \(K>0\), \(\Lambda\ge1\), \(D_{k,w}=D_A\),
and \(A_{\eps,m}=t_0^2/a\). Then
\begin{equation}
 \boxed{k_{w,0}(z)>K\quad(0\le z<H),\qquad k_{w,0}(H)=K.}
 \label{mawf:positive}
\end{equation}
When \(\nu\le d<2\nu\), the profile can initially increase
before decreasing to \(K\); positivity does not require
monotone decrease.
\end{proposition}
\begin{proof}
Use \(L,\beta\) from \cref{mawf:rough-variables}, and set
\begin{equation}
 \begin{gathered}
 G_\beta(x)=\ln(1+x)-\frac{\beta x}{1+x},\qquad
 c=L e^{1-G_\beta(L)},\\
 \Phi(x;c)=
 \begin{cases}
 x/c-G_\beta(x),&x\le c,\\
 1+\ln(x/c)-G_\beta(x),&x\ge c.
 \end{cases}
 \end{gathered}
\end{equation}
Here \(Q_0=A_{\eps,m}\Phi\), \(\Phi(0)=\Phi(L)=0\), and
\begin{equation}
 k_{w,0}(z)-K=\frac{A_{\eps,m}}a[F(az/d)-F(L)],
 \qquad F(x)=\int_0^x\frac{\Phi(t;c)}{1+t}\,\dd t.
 \label{mawf:positivity-F}
\end{equation}
On the tail,
\(\Phi'=[1+(1+\beta)x]/[x(1+x)^2]>0\), so \(\Phi<0\)
before its endpoint zero. On the plateau,
\[
 \Phi'=\frac1c-p_\beta(x),\qquad
 p_\beta(x)=\frac{1+x-\beta}{(1+x)^2},\qquad
 p_\beta'(x)=\frac{2\beta-1-x}{(1+x)^3}.
\]
Thus \(p_\beta\) decreases or has only one maximum.
The allowed signs of \(\Phi\) are negative throughout the
interior, or positive initially and then negative.
This also holds when \(c=L\), since \(\Phi'(L)>0\).
Consequently \(F\) has no interior minimum.

For fixed \(\beta\), write \(I_\beta(L)=F(L)\) while accounting
for \(c=c(L)\). Since
\[
 \frac{c'(L)}{c(L)}
  =\frac{1+(1+\beta)L}{L(1+L)^2}>0,
\]
and \(\Phi(L)=0\), differentiation gives
\begin{equation}
 I_\beta'(L)=-\frac{c'(L)}{c(L)}
 \left[\frac{c-\ln(1+c)}c+\ln\frac{1+L}{1+c}\right]<0.
\end{equation}
Let \(L_1\) be the unique \(G_\beta(L_1)=1\).
At that lower admissible endpoint \(c=L_1\), and
\(T=\ln(1+L_1)\) satisfies \(1<T<2\). Direct integration gives
\begin{equation}
 I_\beta(L_1)
 =(2-T)\left(\frac{2-T}{2}-\frac T{L_1}\right)<0.
\end{equation}
The sign follows from
\(\ln(1+L_1)>2L_1/(L_1+2)\); the derivative of this difference
is \(L_1^2/[(1+L_1)(L_1+2)^2]>0\).
Hence \(F(L)<0=F(0)\) for every admissible \(L\).
The shape argument now gives \(F(x)>F(L)\) for \(x<L\),
which proves \cref{mawf:positive}.
\end{proof}
The proof covers \(d=\nu,\Lambda=1\) algebraically.
Whether such an endpoint is attainable with the original
roughness correlation and \(h_r<H\) is a separate geometric
admissibility check. The stronger proof is specific to the
flat measure and \(D_{k,w}=D_A\); it is not automatically valid
for another scalar diffusivity or a metric-weighted column.

\paragraph{Why the cutoff gate is not a general positivity test.}
On the smooth branch \(h_r=0\), set \(b=10.7\),
\(x=az/\nu\), \(v=10.7\alpha_m\), and
\(M=aH/\nu\). Let \(y_4=146.409759514\ldots\) solve
\[
 \ln(1+y_4)-\frac{y_4}{1+y_4}=4,
\]
and choose \(M=v+y_4\). Then \(\Lambda=4\) and
\(c=az_c/\nu=M e^{-3}\). With
\[
 B(x)=
 \begin{cases}
 0,&x\le v,\\
 \ln(1+x-v)-(x-v)/(1+x-v),&x>v,
 \end{cases}
 \qquad
 \widehat R_c(x)=
 \begin{cases}x/c,&x\le c,\\1+\ln(x/c),&x\ge c,\end{cases}
\]
the exact ground value is
\begin{equation}
 \frac{k_{w,0}(0)}K
 =1-\frac{1}{\kappa^2q_m}
   \int_0^M\frac{\widehat R_c(x)-B(x)}{1+(x-v)_+}\,\dd x.
 \label{mawf:smooth-counterexample}
\end{equation}
At common \(\kappa=0.41\), this gives
\begin{center}
\begin{tabular}{@{}lrrr@{}}
\toprule
Model&A93&A95&E94/E95\\
\midrule
\(k_{w,0}(0)/K\)&\(-0.371380517\)&\(-0.080424984\)&\(-0.436121275\)\\
\bottomrule
\end{tabular}
\end{center}
The budget is admissible, but the stationary kinetic energy is
inadmissible. Thus outside the rough positivity theorem,
the computed profile needs a separate positivity check.

\subsection{Conservative matching and normal pressure balance}
\label{mawf:matching}
At the flat reference,
\begin{equation}
 U_{w,0}(0)=0,\qquad U_{w,0}(H)=U_{o,0}(H),\qquad
 D_A(H)U_{w,0}'(H)=\nu_{t,o,0}(H)U_{o,0}'(H)=t_0.
 \label{mawf:traction}
\end{equation}
Viscosities and gradients need not be continuous.
The inner Reynolds shear is
\(-t_0\nu_{t,A}/D_A\), while outer Reynolds shear is \(-t_0\);
the inner molecular contribution completes total-traction balance.

Use local transported \(k_w\) in the isotropic inner stress.
With \(\Smat_w=(\nabla\bm U_w+\nabla\bm U_w^T)/2\) and
hydrostatic-reduced pressure \(\pi=p/\rho+gz\), the inner
covariance and total stress are
\begin{equation}
 \R_w=\frac23k_w\Id-2\nu_{t,A}\Smat_w,\qquad
 \bm T_w=-\left(\pi_w+\frac23k_w\right)\Id+2D_A\Smat_w .
 \label{mawf:stress}
\end{equation}
The steady vertical balance and its solution are
\begin{equation}
 \left(\pi_{w,0}+\frac23k_{w,0}\right)'=0,\qquad
 \boxed{\pi_{w,0}(z)=\pi_{o,0}+R_{zz,o,0}-\frac23k_{w,0}(z).}
 \label{mawf:pressure}
\end{equation}
Here \(\pi_{o,0}\) and \(R_{zz,o,0}\) are constant in the
outer reference. Define
\begin{equation}
 F_{93}=\frac{q_{93}^2}{1000+q_{93}^3},\qquad
 B_{95}=
 \frac{q_{95}^2\sqrt{1-9/(2q_{95}^2)}}{1+3q_{95}^2}.
\end{equation}
The original outer stresses give
\begin{equation}
 \boxed{\frac{\pi_{w,0}(H)-\pi_{o,0}}K
 =\begin{cases}
 -17F_{93}/3=-0.0630736219,&\text{A93},\\
 -B_{95}=-0.2710910066,&\text{A95},\\
 0,&\text{E94/E95}.
 \end{cases}}
 \label{mawf:pressure-offset}
\end{equation}
Define \(\Delta\pi_H=\pi_{w,0}(H)-\pi_{o,0}\).
Then \(\pi_{w,0}(z)=\pi_{o,0}+\Delta\pi_H
-\frac23[k_{w,0}(z)-K]\).
Even E94/E95 has a varying inner reduced mechanical pressure,
despite its zero interface offset. A pressure variable that already
includes isotropic turbulent stress must be translated consistently.

The \(k\) value and diffusive flux match:
\begin{equation}
 k_w(H)=k_o(H)=K,\qquad
 D_{k,w}(H)k_{w,0}'(H)
 =\frac{\nu_{t,o,0}(H)}{\sigma_k}k_{o,0}'(H)=0.
 \label{mawf:k-interface}
\end{equation}
For the evolving coupled system these are continuity and conservation
conditions, not independent Dirichlet and Neumann prescriptions on
an isolated outer equation. The outer dissipation equation receives
only
\begin{equation}
 \eps_o(H)=E_w(H)=\frac{A_{\eps,m}}H .
 \label{mawf:eps-interface}
\end{equation}
There is no additional two-sided dissipation-flux condition,
because no inner dissipation equation is solved.
When \(z_c=H\), the inner dissipation derivative is zero and the
outer derivative is \(-A_{\eps,m}/H^2\); their traces still agree.

\subsection{Distinction from the constant-inner-\texorpdfstring{\(k\)}{k} EP-AWF}
\label{mawf:comparison}
At identical sampled \(K\), \(h_r,H,\nu\), compatible coefficients
and cutoff, the constant-inner-\(k\) construction uses
\[
 E_w=\frac{A_{\eps,m}}{\max(z,z_c)},\qquad
 C_k^w=E_w-p_0,\qquad S_{k,w}^{\rm EP}=P_A-p_0.
\]
The modified construction uses
\[
 E_w=\frac{A_{\eps,m}}{\max(z,z_c)},\qquad
 C_k^w=0,\qquad S_{k,w}=P_A-E_w.
\]
Both sources have the same layer integral, but their local
balances differ. At the reference, the first uses prescribed
redistribution to support constant \(k=K\); the second uses
ordinary diffusion of \cref{mawf:k-profile}.
The physical reference transport \(-Q_0\) is the same,
while the evolving models are different.

Setting the host's old switch to \(\chi=0\) does not select this
variant when that switch also restores the original-cutoff
dissipation. The present variant retains the budget cutoff
while removing both the prescribed redistribution source and
its equivalent auxiliary flux. Its scalar policy has no
spatial \(\chi\) fade. The earlier bounded-redistribution
positivity remedy would be a further model change if retained.

\subsection{A specified extension to curved terrain}
\label{mawf:curved}
Use the conservative curved momentum completion with the new
endpoint sample \(K=k|_{\Gamma_H}\), compatible coefficients,
metric-weighted reference budget and \(C_k^w=0\).
The equations are restated here so that the scalar replacement
and roughness-force prescription are explicit.

Let \(\bm n\) point from the ground into the fluid and use the
normal map
\begin{equation}
 \bm x(\bm\xi,r)=\bm X(\bm\xi)+r\bm n(\bm\xi),\qquad
 P=\Id-\bm n\otimes\bm n,\quad B_g=\nabla_\Gamma\bm n,\quad
 A_g=I_T+rB_g,\quad J=\det_TA_g.
 \label{mawf:map}
\end{equation}
The inverses of \(A_g\) act tangentially and are extended by zero
normally. The principal stretches \(1+r\gamma_i\) must be
positive, and the normal map must be globally one-to-one on
\(0\le r\le H(\bm\xi)\). The volume measure is
\(J\,\dd S_\Gamma\,\dd r\). Surface derivatives below are
taken at fixed \(r\) before restriction to the interface.

\paragraph{Reference cutoff in the volume measure.}
Set
\begin{equation}
 E_w(r)=\frac{A_{\eps,m}(K)}{\max(r,r_c)},\qquad
 p_0(r;K)=t_0(K)^2\frac{\nu_{t,A}(r;K)}{D_A(r;K)^2}.
\end{equation}
Use counterfactual zero-forcing production in
\begin{equation}
 \boxed{\int_0^H J(r)E_w(r)\,\dd r
        =\mathcal P_J:=\int_0^H J(r)p_0(r;K)\,\dd r.}
 \label{mawf:metric-budget}
\end{equation}
For \(\mathcal B(c)=A_{\eps,m}\int_0^H J/\max(r,c)\,\dd r\),
\[
 \mathcal B'(c)=-\frac{A_{\eps,m}}{c^2}\int_0^cJ\,\dd r<0,
 \qquad \mathcal B(c)\longrightarrow+\infty\quad(c\downarrow0).
\]
Hence a unique \(0<r_c\le H\) exists exactly when
\begin{equation}
 \boxed{\mathcal P_J\ge
        \frac{A_{\eps,m}}H\int_0^H J\,\dd r.}
 \label{mawf:metric-gate}
\end{equation}
At a fixed footpoint, \(J=(1+\gamma_1r)(1+\gamma_2r)\), and
\begin{equation}
 \frac{\mathcal B(c)}{A_{\eps,m}}
 =1+\ln\frac Hc
 +(\gamma_1+\gamma_2)\left(H-\frac c2\right)
 +\gamma_1\gamma_2\left(\frac{H^2}{2}-\frac{c^2}{6}\right).
 \label{mawf:metric-explicit}
\end{equation}
One monotone scalar root solve replaces the flat exponential.

\paragraph{Actual continuity, momentum, forcing and scalar transport.}
Write \(\bm U_w=\bm u_t+w\bm n\) and
\(\bm v=A_g^{-1}\bm u_t\). With Cartesian velocity components,
\begin{equation}
 \partial_r(Jw)+\operatorname{div}_\Gamma(J\bm v)=0,
 \label{mawf:continuity}
\end{equation}
\begin{equation}
 \begin{split}
 J\partial_t\bm U_w
 &+\partial_r\{J(w\bm U_w-T_w\bm n)\}\\
 &+\operatorname{Div}_\Gamma
   \{J(\bm U_w\otimes\bm v-T_wA_g^{-T})\}=J\bm f_w .
 \end{split}
 \label{mawf:curved-momentum}
\end{equation}
Here \(T_w\) is \cref{mawf:stress};
the outer stress is \(T_o=-\pi_o\Id-\R_m\).
The operator \(\operatorname{Div}_\Gamma\) takes the surface
divergence of each Cartesian row. These fluxes already contain
the geometric effects of stress and pressure.

Retain the explicit roughness-force completion
\begin{equation}
 \begin{split}
 a_r(r)&=\mathbf1_{\{r>h_r\}},\qquad
 \bm f_w=\bm f_\pi+(1-a_r)\bm F_A,\\
 \bm F_A&=\mathcal E_H^+P
 \left[\operatorname{Div}_r(\bm u_{o,t}\otimes\bm u_{o,t})
             +\nabla_r\pi_o-\bm f_\pi\right],
 \qquad \bm u_{o,t}=P\bm U_o .
 \end{split}
 \label{mawf:rough-force}
\end{equation}
Here \(\bm f_\pi\) is the remaining body acceleration after the
same hydrostatic-pressure convention used outside,
\(\nabla_rf=A_g^{-T}\nabla_\Gamma f\), and
\[
 \operatorname{Div}_r Q
 =J^{-1}\operatorname{Div}_\Gamma(JQA_g^{-T})
 \quad\text{for }Q\bm n=0.
\]
The sampler \(\mathcal E_H^+\) takes one-sided outer values after
derivatives at fixed \(r\) have been evaluated; extend the resulting
\(\bm F_A\) constantly along each normal ray.
Its work has no guaranteed drag sign and must be retained in
energy accounting.

The flat scalar reconstruction is the declared quasi-steady
parallel-flow reduction of this completion: inner storage,
normal advection, lateral viscous transport and the relevant
\(\nabla_\parallel w\) contributions are neglected, and the
remaining convection/pressure/body-force balance is replaced
by the height-constant sample \(\bm F_A\).
This also replaces the inner effective-pressure gradient by the
sampled outer mechanical-pressure gradient; the discarded
\(\nabla_\parallel(\pi_w+2k_w/3-\pi_o)\) is generally nonzero.
Under those approximations
\((D_A\bm U_\parallel')'=a_r\bm F_A\).
For the curved or unsteady patch, use the full equation
\cref{mawf:curved-momentum}, without additionally imposing that
scalar reduction. All discarded terms vanish on the homogeneous
zero-forcing reference.

The modified actual inner \(k\) equation is
\begin{equation}
 \boxed{\begin{aligned}
 J\partial_t k_w
 &+\partial_r\{J(wk_w-D_{k,w}\partial_rk_w)\}\\
 &+\operatorname{div}_\Gamma
 \{J[\bm v k_w-D_{k,w}A_g^{-1}A_g^{-T}\nabla_\Gamma k_w]\}
 =J(P_A-E_w).
 \end{aligned}}
 \label{mawf:curved-k}
\end{equation}
For incompressible flow,
\[
 P_A=2\nu_{t,A}\Smat_w:\Smat_w .
\]
The isotropic stress uses local \(k_w\), while the viscosity
uses sampled \(K\). If incompressibility is enforced only weakly,
the production contraction \(-\R_w:\nabla\bm U_w\) also contains
\(-(2k_w/3)\nabla\cdot\bm U_w\); its treatment must agree with
the discrete stress-work and pressure convention.
The base-normal tangential traction is
\begin{equation}
 \bm t_A=D_A\left[
 \partial_r\bm u_t+A_g^{-T}\nabla_\Gamma w-B_gA_g^{-1}\bm u_t
 \right].
 \label{mawf:curved-traction}
\end{equation}
Curvature, normal motion and normal pressure balance are thus
retained in the full patch.

The actual source budget remains
\[
 \int_0^HJ(P_A-E_w)\,\dd r
 =\int_0^HJ(P_A-p_0)\,\dd r,
\]
which need not vanish. A frozen reference column can be defined by
\begin{equation}
 Q_J(r)=\int_0^rJ(s)[E_w(s)-p_0(s)]\,\dd s,\qquad
 k_{\rm col}(r)=K-\int_r^H
        \frac{Q_J(s)}{J(s)D_{k,w}(s)}\,\dd s.
 \label{mawf:curved-column}
\end{equation}
Its base-normal endpoint fluxes vanish.
The one-crossing proof for \(d\ge2\nu\) extends to this column
for positive \(J,D_{k,w}\).
The stronger whole-rough-branch proof is not automatically
metric weighted. The column balances only the reference source
and normal diffusion; it is not generally an equilibrium of
the complete terrain-flow equations.

\paragraph{Actual interface geometry and shared fluxes.}
For variable \(H\), put
\begin{equation}
 \bm q_H=A_{g,H}^{-T}\nabla_\Gamma H,\qquad
 \beta_H=\sqrt{1+|\bm q_H|^2},\qquad
 \bm N_H=\frac{\bm n-\bm q_H}{\beta_H}.
 \label{mawf:interface-normal}
\end{equation}
Use common velocity and total traction,
\(\bm U_w=\bm U_o=\widehat{\bm U}\) and
\(T_w\bm N_H=T_o\bm N_H\), with one conservative momentum
flux and opposite orientations on the two sides.
The mass flux per unit ground area is
\(m_H=J_H(w-\bm u_t\cdot\bm q_H)\).
At a common \(k\) trace \(\widehat k\), match
\begin{equation}
 J_H\beta_H(\bm U k-D_k\nabla k)\cdot\bm N_H
 =m_H\widehat k-J_HD_k
 \left[\partial_rk-\bm q_H\cdot A_{g,H}^{-T}\nabla_\Gamma k\right].
 \label{mawf:curved-k-flux}
\end{equation}
Zero base-normal column flux alone does not establish zero flux
through a tilted interface. Surface differentiation of the trace
obeys
\(\nabla_\Gamma K=(\nabla_\Gamma k_w)|_H+
(\partial_rk_w)|_H\nabla_\Gamma H\).
The outer dissipation trace is \(E_w(H)=A_{\eps,m}/H\).
Full normal-traction matching replaces the flat pressure offset;
there is still no independent inner dissipation-flux prescription.

\subsection{Reference realizability, positivity, and stability limits}
\label{mawf:stability}
On a positive rough reference the smallest inner covariance
eigenvalue is
\begin{equation}
 \lambda_{\min}(\R_{w,0})
 =\frac23k_{w,0}-t_0\frac{\nu_{t,A}}{D_A}
 >K\left(\frac23-\frac1{q_m}\right)>0.
 \label{mawf:covariance}
\end{equation}
All three \(q_m\) exceed \(3/2\). Their original outer reference
covariances are also positive definite. These are reference-state
results, not global guarantees for forced or separated states.

With endpoint \(K\), coefficients, cutoff and reference production
held fixed, a wall-normal perturbation \(v=k_w-k_{w,0}\) obeys
\[
 v_t=(D_{k,w}v')',\qquad v(H)=0,\qquad D_{k,w}v'(0)=0,
\]
and
\begin{equation}
 \frac12\frac{\dd}{\dd t}\int_0^H v^2\,\dd z
 =-\int_0^H D_{k,w}(v')^2\,\dd z\le0.
 \label{mawf:scalar-stability}
\end{equation}
This conditional scalar estimate is separate from full coupled
stability.

Dynamic nonnegativity of actual \(k_w\) is not automatic:
at a flat interior zero minimum, positive endpoint \(K\), and
locally zero production and diffusion,
\[
 \partial_t k_w=-E_w(K)<0.
\]
Thus the positive rough reference does not make the evolving
sampled-\(K\) sink positivity preserving. An additional source
modification or numerical positivity device must be identified
and checked separately if introduced.

The wall changes do not alter the original outer constitutive
domains or incremental momentum response. A93 can lose covariance
realizability away from the reference, A95 retains its real-radicand
restriction, and positive covariance in E94/E95 does not ensure
positive incremental momentum diffusion. The previously established
outer constitutive restrictions therefore remain relevant.

\subsection{Illustrative references and implementation implications}
\label{mawf:verification}
Take
\begin{equation}
 \ust=0.4\ \mathrm{m\,s^{-1}},\quad
 \nu=1.5\times10^{-5}\ \mathrm{m^2\,s^{-1}},\quad
 h_r=0.1\ \mathrm m,\quad H=2\ \mathrm m,\quad z_K=z_m=H.
 \label{mawf:inputs}
\end{equation}
Using the original roughness correlation gives
\cref{mawf:numerical-table}. The same sand-grain input is retained,
so aerodynamic roughness is an output rather than a common
prescribed input.

\begin{table}[htbp]
\centering\small
\caption{Analytical rough references. \(K,k(0)\) are in
\(\mathrm{m^2\,s^{-2}}\), and lengths are in metres.}
\label{mawf:numerical-table}
\begin{tabular}{@{}llrrrrr@{}}
\toprule
Model&\(\kappa\) choice&\(K\)&\(k(0)\)&\(\Lambda\)&\(z_c\)&\(z_{0,m}\)\\
\midrule
A93&original&0.54420247&1.08630845&5.90564847&0.01480928&0.00536024\\
A95&original&0.61152007&1.22075495&5.89870963&0.01491240&0.00539286\\
E94/E95&original&0.53192103&1.11770080&5.90535573&0.01481362&0.00535834\\
\midrule
A93&common \(0.41\)&0.54420247&1.13595953&5.90420900&0.01483061&0.00536413\\
A95&common \(0.41\)&0.61152007&1.20319430&5.89920134&0.01490507&0.00539152\\
E94/E95&common \(0.41\)&0.53192103&1.12369422&5.90518519&0.01481614&0.00535880\\
\bottomrule
\end{tabular}
\end{table}

All six cases satisfy \(d\ge2\nu\) and \(\Lambda>1\), so their
positive monotone profiles have an analytical guarantee.
The interface pressure offsets are
\(-0.034324820606\), \(-0.165777590815\), and \(0\)
\(\mathrm{m^2\,s^{-2}}\), respectively.
At fixed \(\ust\) these offsets do not change with \(\kappa\),
because \(K=q_m\ust^2\) is unchanged; the inner pressure profile
does change with \(k_w(z)\).

Independent quadrature of the production and dissipation budgets,
the cumulative defect, and the closed primitive agrees with the
analytical expressions. Separately integrated source and diffusion
moments for Legendre degrees \(0,\ldots,6\) give a maximum absolute
weak residual of \(2.33\times10^{-16}\) across these cases;
closed \(k\) versus quadrature differs by at most
\(6.89\times10^{-15}\) in SI units. These are checks of the
formulas, not CFD or terrain validation.

The model-specific \(\kappa\) affects the viscosity slope,
dissipation amplitude, cutoff and inner profile. A further useful
identity is
\[
 \frac{\kappa^2[k_{w,0}(z)-K]}{u_m^2}
 =\mathcal F(az/d)-\mathcal F(aH/d).
\]
At common \(\kappa,\ust,H,z_0,\nu\), if each roughness input
is calibrated to the same \(z_0\), \cref{mawf:d-calibration}
makes \(a,d\) identical. Then the zero-forcing wall velocity,
production, sink and cumulative defect coincide across the models,
and their \(k\) profiles differ only by \(K=q_m\ust^2\).
Their forced responses can still differ through the calibrated
\(h_r\) in the forcing cutoff.

For a spectral/\(hp\) implementation, the exact continuum weak
identity \cref{mawf:weak} does not by itself make a finite polynomial
interpolant stationary. The ordinary diffusive flux, nonzero source
moments and shared interface terms must be represented and integrated
consistently. A local discrete wall solve using the same source,
quadrature and boundary contract can supply a stationary discrete
profile. Checking only the cell-mean energy budget is insufficient.

The selected variant therefore consists of endpoint sampling,
the two compatible coefficients, the budget cutoff and its plateau,
retained finite-\(\nu\) transport with \(P_A-E_w\), and conservative
velocity/stress/scalar coupling. Its exactness claim is the
flat neutral reference of these specified equations. Positivity
outside the rough theorem, global coupled stability and fully
discrete preservation remain separate requirements.

% The main note already calls \appendix; do not call it again here.
% Uses the main note's amsmath/amssymb/amsthm/bm/booktabs/natbib/cleveref
% packages, proposition environment, and \eps, \ust, \dd, \Id, \R, \Smat.
% Bibliography keys: Shih1993, Shih1995ASM, Shih1994EVM, Shih1995EVM,
% and SCI2006. All labels introduced here have the unique prefix mawf:.
% No document class, document environment, or bibliography is added.

\section{AWF with coefficients compatible with A93, A95 and E94/E95,
a budget-determined cutoff, and a nonconstant inner kinetic-energy profile}
\label{mawf:appendix}

This appendix specifies a modified analytical wall function for each of
the three outer closures A93, A95 and E94/E95
\citep{Shih1993,Shih1995ASM,Shih1994EVM,Shih1995EVM}.
The objective is to retain its exact neutral outer equilibrium while
allowing kinetic energy to vary inside the modeled wall layer.
Two auxiliary coefficients are matched to the selected outer branch;
the dissipation cutoff is determined by the reference energy budget;
and the prescribed redistribution source is removed. Ordinary diffusion
of the transported inner kinetic energy then balances the local
production--dissipation difference.

The rough-wall viscosity kernel and analytical momentum reconstruction
start from the AWF of Suga, Craft and Iacovides \citep{SCI2006}.
The compatible coefficients, endpoint sampling, budget cutoff and
variable inner reference specified below are explicit modifications.
Selecting this variant replaces the earlier wall-coefficient,
sampling and \(\chi\)-dependent scalar-source choices. Its own
dissipation law is retained with \(C_k^w=0\).

\subsection{Scope, endpoint sampling, and compatible coefficients}
\label{mawf:coefficients}

Use the selected high-Re outer equations: molecular momentum and scalar
diffusion are omitted from the outer analytical reference, and the
E94/E95 dissipation denominator is reduced from
\(k+\sqrt{\nu\eps}\) to \(k\) in that outer limit.
The inner AWF retains the fixed physical \(\nu>0\) in its momentum,
production and scalar-diffusion formulas. These are two declared
equation regimes; the result is not an exact logarithmic solution of
an unchanged finite-Re outer system.

Only inner \(k_w\) is transported. Inner dissipation is a dependent
total sink, and no inner dissipation transport equation is introduced.
The inner stress uses the auxiliary linear viscosity specified below;
the original A93/A95 nonlinear stress laws remain the outer closures.
Coordinates are \(x\) streamwise, \(y\) spanwise, and \(z\) vertical.
All stresses are per unit density.

Let \(m\in\{93,95,E\}\) identify the outer model.
Its neutral equilibrium ratio \(q_m=k_{o,0}/\ust^2\) is
\begin{equation}
 \boxed{
 q_{93}=\frac{57+\sqrt{6249}}{40},\qquad
 q_{95}=\frac{3+\sqrt{61}}{2\sqrt2},\qquad
 q_E=\frac{3+\sqrt{41}}{2\sqrt2}.}
 \label{mawf:q}
\end{equation}
These roots follow from \(C_{\mu,m}(q_m)q_m^2=1\) in simple shear.
The nonlinear corrections of A93 and A95 have zero \(xz\) component
there, but the full Reynolds shear is \(R_{xz,o,0}=-\ust^2\).

The outer dissipation equation additionally requires
\begin{equation}
 \boxed{\kappa^2q_m=\sigma_{\eps,m}\Delta_m,\qquad
 \Delta_{93}=\Delta_{95}=0.48,\qquad
 \Delta_E=1.9-0.43q_E.}
 \label{mawf:kappa}
\end{equation}
For E94/E95 the branch \(C_1=0.43\) is active at \(q_E\).
Retaining the original outer coefficients means selecting
\(\kappa_m=\sqrt{\sigma_{\eps,m}^{\rm orig}\Delta_m/q_m}\).
A prescribed common \(\kappa\) instead requires the single outer
replacement \(\sigma_{\eps,m}^{\rm eq}=\kappa^2q_m/\Delta_m\).
The exact values implied by the printed coefficients are listed in
\cref{mawf:outer-table}; they need not coincide with a nominal
\(\kappa\) used in the original empirical calibration.

\begin{table}[htbp]
\centering\small
\caption{Outer choices. The last column is the required one-coefficient
change when a common \(\kappa=0.41\) is selected.}
\label{mawf:outer-table}
\begin{tabular}{@{}lrrrr@{}}
\toprule
Model & \(q_m\) & original \(\sigma_\eps\)& compatible \(\kappa_m\)&
\(\sigma_\eps^{\,0.41}\)\\
\midrule
A93 &3.401265417&1.3&0.428323629&1.191151493\\
A95 &3.822000426&1.3&0.404060995&1.338496399\\
E94/E95 &3.324506456&1.2&0.412087472&1.187873345\\
\bottomrule
\end{tabular}
\end{table}

Use endpoint sampling and tail matching:
\begin{equation}
 \boxed{z_K=z_m=H,\qquad K=k_w(H)=k_o(H)>0.}
 \label{mawf:sampling}
\end{equation}
The subscripts \(w,o\) denote wall layer and outer region; subscript
\(0\) denotes their steady flat reference. At each positive current
sample \(K\), define counterfactual zero-forcing scales
\begin{equation}
 u_m(K)=\sqrt{\frac K{q_m}},\qquad
 t_0(K)=\frac K{q_m},\qquad
 A_{\eps,m}(K)=\frac{u_m(K)^3}{\kappa}.
 \label{mawf:scales}
\end{equation}
At the chosen inlet reference \(K=q_m\ust^2\), \(u_m=\ust\).
Away from that reference, \(t_0\) is not imposed as the actual
wall traction. Model and sample arguments are suppressed below.

Match the wall dissipation-tail amplitude and viscosity slope:
\begin{equation}
 \frac{K^{3/2}}{c_{\ell,m}^w}=\frac{u_m^3}{\kappa},\qquad
 c_{\mu,m}^wc_{\ell,m}^w\sqrt K=\kappa u_m.
\end{equation}
The two compatible auxiliary coefficients and their product are
\begin{equation}
 \boxed{c_{\ell,m}^w=\kappa q_m^{3/2},\qquad
 c_{\mu,m}^w=q_m^{-2},\qquad
 \alpha_m=c_{\mu,m}^wc_{\ell,m}^w=\frac{\kappa}{\sqrt{q_m}}.}
 \label{mawf:wall-coefficients}
\end{equation}
Consequently, with \(a=\alpha_m\sqrt K\),
\begin{equation}
 \boxed{a=\kappa u_m,\qquad
 \frac{K^{3/2}}{c_{\ell,m}^w}=A_{\eps,m}=\frac{t_0^2}{a}.}
 \label{mawf:identities}
\end{equation}
These two matching requirements uniquely determine the two auxiliary
coefficients. Traction continuity by itself does not require
equal inner and outer viscosity slopes.
The auxiliary \(c_\mu^w\) is constant for the selected branch;
it is not the current gradient-dependent outer \(C_\mu\).

\begin{table}[htbp]
\centering\small
\caption{Compatible wall coefficients. Each row uses the same
\(\kappa\) in the outer equations and the wall model.}
\label{mawf:coefficient-table}
\begin{tabular}{@{}llrrrr@{}}
\toprule
Model&outer choice&\(\kappa\)&\(c_\mu^w\)&\(c_\ell^w\)&\(\alpha\)\\
\midrule
A93&original&0.428323629&0.086440835&2.686784401&0.232247887\\
A95&original&0.404060995&0.068457106&3.019138428&0.206681480\\
E94/E95&original&0.412087472&0.090478551&2.497929531&0.226009045\\
\midrule
A93&common \(\kappa\)&0.410000000&0.086440835&2.571844117&0.222312353\\
A95&common \(\kappa\)&0.410000000&0.068457106&3.063514593&0.209719344\\
E94/E95&common \(\kappa\)&0.410000000&0.090478551&2.485276011&0.224864173\\
\bottomrule
\end{tabular}
\end{table}

Retain the original roughness-intercept correlation. With sand-grain
roughness \(h_r\ge0\), define, for \(h_r>0\),
\begin{equation}
 h^*=\frac{h_r\sqrt K}{\nu},\quad x=\frac{h^*}{70},\quad
 m_r=\max(0.5-0.4x^{0.7},\,1-0.79x^{-0.28}),\quad
 b=10.7(1-x^{m_r});
 \label{mawf:roughness}
\end{equation}
for \(h_r=0\), set \(b=10.7\). The wall coefficients are
\begin{equation}
 \begin{gathered}
 z_v=\frac{\nu b}{\sqrt K},\qquad
 \nu_{t,A}(z)=a(z-z_v)_+,\qquad D_A(z)=\nu+\nu_{t,A}(z),\\
 D_{k,w}(z)=\nu+\frac{\nu_{t,A}(z)}{\sigma_k}=D_A(z).
 \end{gathered}
 \label{mawf:viscosity}
\end{equation}
Here \((r)_+=\max(r,0)\), and the last equality uses \(\sigma_k=1\)
for all three models. These coefficients sample endpoint \(K\);
local transported \(k_w(z)\) is allowed to vary. Substituting local
\(k_w\) into the kernel would define a different model. Likewise,
an interior sample or average cannot simply replace the endpoint
sample once the inner reference is nonconstant.

\subsection{Momentum, production, and the outer velocity anchor}
\label{mawf:momentum}
For a flat quasi-steady column, put
\begin{equation}
 L_h(z)=(z-h_r)_+,\quad
 J_0(z)=\int_0^z\frac{\dd s}{D_A(s)},\quad
 J_1(z)=\int_0^z\frac{L_h(s)}{D_A(s)}\,\dd s,\quad
 J_2(z)=\int_0^z\frac{\dd s}{D_A(s)^2}.
 \label{mawf:J}
\end{equation}
With prescribed constant tangential forcing \(F\), retain
\begin{equation}
 (D_AU_A')'=F\,\mathbf1_{\{z>h_r\}},\qquad
 U_A(0)=0,\quad U_A(H)=U_H.
 \label{mawf:momentum-bvp}
\end{equation}
Integration yields the original analytical form, now evaluated
with the compatible viscosity:
\begin{equation}
 \boxed{\begin{aligned}
 t_A&=\frac{U_H-FJ_1(H)}{J_0(H)},&
 U_A(z)&=t_AJ_0(z)+FJ_1(z),\\
 P_A(z)&=\nu_{t,A}(z)
       \left|\frac{t_A+FL_h(z)}{D_A(z)}\right|^2 .
 \end{aligned}}
 \label{mawf:forced}
\end{equation}
The upper traction is \(T_H=t_A+FL_h(H)\), generally different
from ground traction \(t_A\). The vector tangential form uses
componentwise reconstruction and the Euclidean square in \(P_A\).
The coordinates remain the physical Cartesian coordinates.

At the counterfactual zero-forcing reference choose
\begin{equation}
 U_{H,0}=t_0J_0(H),\qquad U_{w,0}(z)=t_0J_0(z),\qquad
 p_0(z)=t_0^2\frac{\nu_{t,A}(z)}{D_A(z)^2}.
 \label{mawf:reference-production}
\end{equation}
Its integrated production is exactly
\begin{equation}
 \boxed{\mathcal P_0=\int_0^H p_0\,\dd z
     =t_0^2[J_0(H)-\nu J_2(H)].}
 \label{mawf:production-budget}
\end{equation}
Here \(p_0\) is pointwise reference production and
\(\mathcal P_0/H\) is its height mean. Neither is the actual
production of a general current flow. The molecular part of mean
viscous dissipation is excluded from turbulent production; hence
the term \(-\nu J_2\) must be retained.

For \(z\ge H\), the compatible outer reference is
\begin{equation}
 \boxed{\begin{aligned}
 U_{o,0}(z)&=t_0J_0(H)+\frac{u_m}{\kappa}\ln\frac zH,\\
 k_{o,0}(z)&=K,\qquad
 \eps_{o,0}(z)=\frac{A_{\eps,m}}z,\qquad
 \nu_{t,o,0}(z)=az .
 \end{aligned}}
 \label{mawf:outer-reference}
\end{equation}
Its aerodynamic intercept is
\begin{equation}
 z_{0,m}=H\exp[-\kappa u_mJ_0(H)]=H\exp[-aJ_0(H)].
 \label{mawf:z0}
\end{equation}
Keeping the same \(h_r\) after changing coefficients generally
changes \(z_{0,m}\). A separately prescribed \(z_0\) requires
\begin{equation}
 J_0(H)=\frac{\ln(H/z_0)}{\kappa\ust}
 \label{mawf:z0-calibration}
\end{equation}
at the selected inlet reference. A physical roughness input calibrated
there is subsequently held fixed as the endpoint sample evolves.

\subsection{Budget-determined dissipation and admissibility}
\label{mawf:cutoff-section}
Use a positive continuous dependent sink and no prescribed redistribution:
\begin{equation}
 \boxed{E_w(z)=\frac{A_{\eps,m}}{\max(z,z_c)},\qquad
 C_k^w(z)=0,\qquad 0<z_c\le H.}
 \label{mawf:dissipation}
\end{equation}
Its integral and upper trace are
\begin{equation}
 \int_0^HE_w\,\dd z=A_{\eps,m}\left(1+\ln\frac H{z_c}\right),
 \qquad E_w(H)=\frac{A_{\eps,m}}H.
 \label{mawf:dissipation-integral}
\end{equation}
Zero ground \(k\) flux and the constant outer reference require
\(\int_0^H(E_w-p_0)\,\dd z=0\). Thus
\begin{equation}
 \boxed{\Lambda=\frac{\mathcal P_0}{A_{\eps,m}}
       =a[J_0(H)-\nu J_2(H)],\qquad
 z_c=H\exp(1-\Lambda).}
 \label{mawf:cutoff}
\end{equation}
The integral in \cref{mawf:dissipation-integral} decreases strictly
with \(z_c\), tends to infinity as \(z_c\downarrow0\), and equals
\(A_{\eps,m}\) at \(z_c=H\). Therefore
\begin{equation}
 \boxed{\text{a unique }0<z_c\le H\text{ exists}
        \quad\Longleftrightarrow\quad \Lambda\ge1.}
 \label{mawf:gate}
\end{equation}
If the gate fails, this sink family cannot satisfy both the budget
and the outer trace. Clipping \(z_c\) to \(H\) loses the budget.

The original molecular cutoff
\(z_{\eps,A}=2c_\ell^w\nu/\sqrt K\) is replaced by
\cref{mawf:cutoff}. The plateau is \(A_{\eps,m}/z_c\).
Imposing also \(E_w(0)=2\nu K/z_c^2\) would force the original
molecular cutoff again. The replacement cutoff and plateau are
explicit modifications of the dissipation law.
The retained finite-viscosity contribution also gives
\begin{equation}
 \boxed{z_c=e\,z_{0,m}\exp[a\nu J_2(H)]>e\,z_{0,m},}
 \label{mawf:cutoff-z0}
\end{equation}
and the mechanical-work budget
\begin{equation}
 t_0U_{H,0}
 =\nu t_0^2J_2(H)+\int_0^Hp_0\,\dd z
 =\nu t_0^2J_2(H)+\int_0^HE_w\,\dd z.
 \label{mawf:work}
\end{equation}

\subsection{Nonconstant inner kinetic energy and every weak test mode}
\label{mawf:inner-k}
Retain the actual inner transport equation
\begin{equation}
 \partial_tk_w+\bm U_w\cdot\nabla k_w
 =\nabla\cdot(D_{k,w}\nabla k_w)+P_A-E_w.
 \label{mawf:k-transport}
\end{equation}
The steady parallel reference satisfies
\begin{equation}
 (D_{k,w}k_{w,0}')'=E_w-p_0,\qquad
 D_{k,w}(0)k_{w,0}'(0)=0,\qquad k_{w,0}(H)=K.
 \label{mawf:k-bvp}
\end{equation}
Define
\begin{equation}
 Q_0(z)=\int_0^z[E_w(s)-p_0(s)]\,\dd s.
\end{equation}
Two integrations give the unique reference profile
\begin{equation}
 \boxed{
 D_{k,w}k_{w,0}'=Q_0,\qquad
 k_{w,0}(z)=K-\int_z^H\frac{Q_0(s)}{D_{k,w}(s)}\,\dd s.}
 \label{mawf:k-profile}
\end{equation}
The physical upward diffusive flux is \(-Q_0\).
By definition \(Q_0(0)=0\), and the budget gives \(Q_0(H)=0\).
Hence
\begin{equation}
 k_{w,0}(H)=K,\qquad
 k_{w,0}'(0)=k_{w,0}'(H)=0.
 \label{mawf:k-endpoints}
\end{equation}
The ground value is determined by \cref{mawf:k-profile};
an additional \(k_w(0)=0\) generally overdetermines this reference.
The upper zero flux follows from the cutoff budget rather than
being an additional independent boundary condition on a fixed-cutoff
inner problem.

For any sufficiently regular test function \(\phi\),
\begin{equation}
 \boxed{
 -\int_0^H D_{k,w}k_{w,0}'\phi'\,\dd z
 +\int_0^H(p_0-E_w)\phi\,\dd z=0.}
 \label{mawf:weak}
\end{equation}
Indeed, substitution reduces the left side to
\(-[Q_0\phi]_0^H=0\). The constant test function expresses the
budget; for nonconstant modes, ordinary diffusion cancels nonzero
source moments. This is the mechanism required in a spectral/\(hp\)
wall element.

The analytical profile specifies a reference and an initialization;
the evolving \(k_w\) is transported, not reset to this profile.
Any equivalent prescribed redistribution flux from an earlier
flux-form implementation is also removed. The physical flux
\(-D_{k,w}\nabla k_w\) supplies the reference transport.
The cutoff continues to use counterfactual \(p_0(K)\).
For actual production,
\begin{equation}
 \int_0^H(P_A-E_w)\,\dd z=\int_0^H(P_A-p_0)\,\dd z
 \label{mawf:actual-budget}
\end{equation}
need not vanish. Determining the cutoff from actual \(P_A\)
would define a different model.

\subsection{Closed rough-branch formulas and positivity}
\label{mawf:rough}
Assume \(z_v\le0\), and put
\begin{equation}
 d=\nu-az_v\ge\nu,\qquad D_A=d+az,\qquad
 \zeta=\frac{az}{d},\quad L=\frac{aH}{d},\quad
 \beta=\frac{\nu}{d}\in(0,1].
 \label{mawf:rough-variables}
\end{equation}
Then
\begin{align}
 J_0(z)&=\frac1a\ln\frac{d+az}{d},&
 J_2(z)&=\frac1a\left(\frac1d-\frac1{d+az}\right),\\
 p_0(z)&=t_0^2\frac{d+az-\nu}{(d+az)^2},&
 \int_0^zp_0\,\dd s
 &=A_{\eps,m}\left[\ln\frac{d+az}{d}
          -\nu\left(\frac1d-\frac1{d+az}\right)\right].
 \label{mawf:rough-integrals}
\end{align}
Consequently,
\begin{equation}
 \boxed{\begin{gathered}
 \Lambda=\ln(1+L)-\frac{\beta L}{1+L},\qquad
 z_{0,m}=\frac{dH}{aH+d},\\
 z_c=e\,z_{0,m}
 \exp\left[\nu\left(\frac1d-\frac1{aH+d}\right)\right].
 \end{gathered}}
 \label{mawf:rough-cutoff}
\end{equation}
At a prescribed \(z_0<H\), the required rough-branch intercept is
\begin{equation}
 d=\frac{aH z_0}{H-z_0}.
 \label{mawf:d-calibration}
\end{equation}
It must be attainable by the roughness correlation and the selected
wall-layer geometry, normally \(h_r<H\).
Equal viscosity slopes do not remove the finite intercept:
\(D_A(H)=d+aH\), while the outer high-Re viscosity is \(aH\).

Define
\begin{equation}
 R_c(z)=
 \begin{cases}
 z/z_c,&0\le z\le z_c,\\
 1+\ln(z/z_c),&z_c\le z\le H.
 \end{cases}
\end{equation}
The cumulative defect is
\begin{equation}
 \boxed{Q_0(z)=A_{\eps,m}
 \left[R_c(z)-\ln\frac{d+az}{d}
       +\nu\left(\frac1d-\frac1{d+az}\right)\right].}
 \label{mawf:rough-Q}
\end{equation}

\paragraph{Closed kinetic-energy primitive.}
For \(c=az_c/d=L e^{1-\Lambda}\), set
\begin{equation}
 \mathcal F_<(\zeta)=
 \frac{\zeta-\ln(1+\zeta)}c-\frac12\ln^2(1+\zeta)
 +\beta\left[\ln(1+\zeta)+\frac1{1+\zeta}-1\right],
\end{equation}
and
\begin{equation}
 \begin{split}
 \mathcal G(\zeta)
 &=\left[1+\ln\frac{\zeta}{c}+\beta\right]\ln(1+\zeta)
   +\operatorname{Li}_2(-\zeta)
   -\frac12\ln^2(1+\zeta)+\frac{\beta}{1+\zeta},\\
 \operatorname{Li}_2(-\zeta)
 &=-\int_0^\zeta\frac{\ln(1+s)}s\,\dd s\qquad(\zeta\ge0).
 \end{split}
\end{equation}
Join the primitives:
\begin{equation}
 \mathcal F(\zeta)=
 \begin{cases}
 \mathcal F_<(\zeta),&\zeta\le c,\\
 \mathcal F_<(c)+\mathcal G(\zeta)-\mathcal G(c),&\zeta\ge c.
 \end{cases}
\end{equation}
Since
\(\mathcal F'(\zeta)=Q_0/[A_{\eps,m}(1+\zeta)]\),
\begin{equation}
 \boxed{k_{w,0}(z)=K+\frac{A_{\eps,m}}a
  [\mathcal F(az/d)-\mathcal F(aH/d)],\qquad
  \frac{A_{\eps,m}}a=\frac{u_m^2}{\kappa^2}.}
 \label{mawf:k-closed}
\end{equation}
Definite quadrature of \cref{mawf:k-profile} is an equivalent
evaluation. Near \(H\), it can avoid cancellation of close
primitive values.

\paragraph{A sufficient condition for monotone decrease.}
For \(z\ge z_c>0\), direct simplification using
\(A_{\eps,m}=t_0^2/a\) gives
\begin{equation}
 E_w-p_0
 =\frac{A_{\eps,m}[(d+\nu)az+d^2]}{z(d+az)^2}>0.
 \label{mawf:tail-defect}
\end{equation}
Moreover,
\begin{equation}
 p_0'(z)=\frac{a t_0^2[2\nu-(d+az)]}{(d+az)^3}.
 \label{mawf:production-derivative}
\end{equation}

\begin{proposition}[Monotone positive rough reference]
\label{mawf:monotone-proposition}
If \(d\ge2\nu\), \(K>0\), and \(\Lambda\ge1\), then
\begin{equation}
 Q_0(z)<0,\qquad k_{w,0}'(z)<0,\qquad
 k_{w,0}(z)>K>0\quad(0<z<H),
\end{equation}
and \(k_{w,0}(0)>K\).
\end{proposition}
\begin{proof}
The production decreases strictly for \(z>0\).
On the dissipation plateau, \(E_w-p_0\) therefore increases
strictly, and its value at \(z_c\) is positive by
\cref{mawf:tail-defect}. Its zero total integral forces one
negative-to-positive crossing strictly inside \((0,z_c)\).
The cumulative defect is negative until it returns to zero at
\(H\), and \cref{mawf:k-profile} gives the result.
This argument includes \(z_c=H\): the positive endpoint
defect still forces a crossing below \(H\).
\end{proof}
The one-crossing argument works for any positive \(D_{k,w}\).
The following stronger positivity statement uses specifically
\(\sigma_k=1\), hence \(D_{k,w}=D_A\).

\begin{proposition}[Positive reference on the whole rough branch]
\label{mawf:positive-proposition}
Assume the flat measure, constant physical \(\nu>0\),
\(d\ge\nu\), \(K>0\), \(\Lambda\ge1\), \(D_{k,w}=D_A\),
and \(A_{\eps,m}=t_0^2/a\). Then
\begin{equation}
 \boxed{k_{w,0}(z)>K\quad(0\le z<H),\qquad k_{w,0}(H)=K.}
 \label{mawf:positive}
\end{equation}
When \(\nu\le d<2\nu\), the profile can initially increase
before decreasing to \(K\); positivity does not require
monotone decrease.
\end{proposition}
\begin{proof}
Use \(L,\beta\) from \cref{mawf:rough-variables}, and set
\begin{equation}
 \begin{gathered}
 G_\beta(x)=\ln(1+x)-\frac{\beta x}{1+x},\qquad
 c=L e^{1-G_\beta(L)},\\
 \Phi(x;c)=
 \begin{cases}
 x/c-G_\beta(x),&x\le c,\\
 1+\ln(x/c)-G_\beta(x),&x\ge c.
 \end{cases}
 \end{gathered}
\end{equation}
Here \(Q_0=A_{\eps,m}\Phi\), \(\Phi(0)=\Phi(L)=0\), and
\begin{equation}
 k_{w,0}(z)-K=\frac{A_{\eps,m}}a[F(az/d)-F(L)],
 \qquad F(x)=\int_0^x\frac{\Phi(t;c)}{1+t}\,\dd t.
 \label{mawf:positivity-F}
\end{equation}
On the tail,
\(\Phi'=[1+(1+\beta)x]/[x(1+x)^2]>0\), so \(\Phi<0\)
before its endpoint zero. On the plateau,
\[
 \Phi'=\frac1c-p_\beta(x),\qquad
 p_\beta(x)=\frac{1+x-\beta}{(1+x)^2},\qquad
 p_\beta'(x)=\frac{2\beta-1-x}{(1+x)^3}.
\]
Thus \(p_\beta\) decreases or has only one maximum.
The allowed signs of \(\Phi\) are negative throughout the
interior, or positive initially and then negative.
This also holds when \(c=L\), since \(\Phi'(L)>0\).
Consequently \(F\) has no interior minimum.

For fixed \(\beta\), write \(I_\beta(L)=F(L)\) while accounting
for \(c=c(L)\). Since
\[
 \frac{c'(L)}{c(L)}
  =\frac{1+(1+\beta)L}{L(1+L)^2}>0,
\]
and \(\Phi(L)=0\), differentiation gives
\begin{equation}
 I_\beta'(L)=-\frac{c'(L)}{c(L)}
 \left[\frac{c-\ln(1+c)}c+\ln\frac{1+L}{1+c}\right]<0.
\end{equation}
Let \(L_1\) be the unique \(G_\beta(L_1)=1\).
At that lower admissible endpoint \(c=L_1\), and
\(T=\ln(1+L_1)\) satisfies \(1<T<2\). Direct integration gives
\begin{equation}
 I_\beta(L_1)
 =(2-T)\left(\frac{2-T}{2}-\frac T{L_1}\right)<0.
\end{equation}
The sign follows from
\(\ln(1+L_1)>2L_1/(L_1+2)\); the derivative of this difference
is \(L_1^2/[(1+L_1)(L_1+2)^2]>0\).
Hence \(F(L)<0=F(0)\) for every admissible \(L\).
The shape argument now gives \(F(x)>F(L)\) for \(x<L\),
which proves \cref{mawf:positive}.
\end{proof}
The proof covers \(d=\nu,\Lambda=1\) algebraically.
Whether such an endpoint is attainable with the original
roughness correlation and \(h_r<H\) is a separate geometric
admissibility check. The stronger proof is specific to the
flat measure and \(D_{k,w}=D_A\); it is not automatically valid
for another scalar diffusivity or a metric-weighted column.

\paragraph{Why the cutoff gate is not a general positivity test.}
On the smooth branch \(h_r=0\), set \(b=10.7\),
\(x=az/\nu\), \(v=10.7\alpha_m\), and
\(M=aH/\nu\). Let \(y_4=146.409759514\ldots\) solve
\[
 \ln(1+y_4)-\frac{y_4}{1+y_4}=4,
\]
and choose \(M=v+y_4\). Then \(\Lambda=4\) and
\(c=az_c/\nu=M e^{-3}\). With
\[
 B(x)=
 \begin{cases}
 0,&x\le v,\\
 \ln(1+x-v)-(x-v)/(1+x-v),&x>v,
 \end{cases}
 \qquad
 \widehat R_c(x)=
 \begin{cases}x/c,&x\le c,\\1+\ln(x/c),&x\ge c,\end{cases}
\]
the exact ground value is
\begin{equation}
 \frac{k_{w,0}(0)}K
 =1-\frac{1}{\kappa^2q_m}
   \int_0^M\frac{\widehat R_c(x)-B(x)}{1+(x-v)_+}\,\dd x.
 \label{mawf:smooth-counterexample}
\end{equation}
At common \(\kappa=0.41\), this gives
\begin{center}
\begin{tabular}{@{}lrrr@{}}
\toprule
Model&A93&A95&E94/E95\\
\midrule
\(k_{w,0}(0)/K\)&\(-0.371380517\)&\(-0.080424984\)&\(-0.436121275\)\\
\bottomrule
\end{tabular}
\end{center}
The budget is admissible, but the stationary kinetic energy is
inadmissible. Thus outside the rough positivity theorem,
the computed profile needs a separate positivity check.

\subsection{Conservative matching and normal pressure balance}
\label{mawf:matching}
At the flat reference,
\begin{equation}
 U_{w,0}(0)=0,\qquad U_{w,0}(H)=U_{o,0}(H),\qquad
 D_A(H)U_{w,0}'(H)=\nu_{t,o,0}(H)U_{o,0}'(H)=t_0.
 \label{mawf:traction}
\end{equation}
Viscosities and gradients need not be continuous.
The inner Reynolds shear is
\(-t_0\nu_{t,A}/D_A\), while outer Reynolds shear is \(-t_0\);
the inner molecular contribution completes total-traction balance.

Use local transported \(k_w\) in the isotropic inner stress.
With \(\Smat_w=(\nabla\bm U_w+\nabla\bm U_w^T)/2\) and
hydrostatic-reduced pressure \(\pi=p/\rho+gz\), the inner
covariance and total stress are
\begin{equation}
 \R_w=\frac23k_w\Id-2\nu_{t,A}\Smat_w,\qquad
 \bm T_w=-\left(\pi_w+\frac23k_w\right)\Id+2D_A\Smat_w .
 \label{mawf:stress}
\end{equation}
The steady vertical balance and its solution are
\begin{equation}
 \left(\pi_{w,0}+\frac23k_{w,0}\right)'=0,\qquad
 \boxed{\pi_{w,0}(z)=\pi_{o,0}+R_{zz,o,0}-\frac23k_{w,0}(z).}
 \label{mawf:pressure}
\end{equation}
Here \(\pi_{o,0}\) and \(R_{zz,o,0}\) are constant in the
outer reference. Define
\begin{equation}
 F_{93}=\frac{q_{93}^2}{1000+q_{93}^3},\qquad
 B_{95}=
 \frac{q_{95}^2\sqrt{1-9/(2q_{95}^2)}}{1+3q_{95}^2}.
\end{equation}
The original outer stresses give
\begin{equation}
 \boxed{\frac{\pi_{w,0}(H)-\pi_{o,0}}K
 =\begin{cases}
 -17F_{93}/3=-0.0630736219,&\text{A93},\\
 -B_{95}=-0.2710910066,&\text{A95},\\
 0,&\text{E94/E95}.
 \end{cases}}
 \label{mawf:pressure-offset}
\end{equation}
Define \(\Delta\pi_H=\pi_{w,0}(H)-\pi_{o,0}\).
Then \(\pi_{w,0}(z)=\pi_{o,0}+\Delta\pi_H
-\frac23[k_{w,0}(z)-K]\).
Even E94/E95 has a varying inner reduced mechanical pressure,
despite its zero interface offset. A pressure variable that already
includes isotropic turbulent stress must be translated consistently.

The \(k\) value and diffusive flux match:
\begin{equation}
 k_w(H)=k_o(H)=K,\qquad
 D_{k,w}(H)k_{w,0}'(H)
 =\frac{\nu_{t,o,0}(H)}{\sigma_k}k_{o,0}'(H)=0.
 \label{mawf:k-interface}
\end{equation}
For the evolving coupled system these are continuity and conservation
conditions, not independent Dirichlet and Neumann prescriptions on
an isolated outer equation. The outer dissipation equation receives
only
\begin{equation}
 \eps_o(H)=E_w(H)=\frac{A_{\eps,m}}H .
 \label{mawf:eps-interface}
\end{equation}
There is no additional two-sided dissipation-flux condition,
because no inner dissipation equation is solved.
When \(z_c=H\), the inner dissipation derivative is zero and the
outer derivative is \(-A_{\eps,m}/H^2\); their traces still agree.

\subsection{Distinction from the constant-inner-\texorpdfstring{\(k\)}{k} EP-AWF}
\label{mawf:comparison}
At identical sampled \(K\), \(h_r,H,\nu\), compatible coefficients
and cutoff, the constant-inner-\(k\) construction uses
\[
 E_w=\frac{A_{\eps,m}}{\max(z,z_c)},\qquad
 C_k^w=E_w-p_0,\qquad S_{k,w}^{\rm EP}=P_A-p_0.
\]
The modified construction uses
\[
 E_w=\frac{A_{\eps,m}}{\max(z,z_c)},\qquad
 C_k^w=0,\qquad S_{k,w}=P_A-E_w.
\]
Both sources have the same layer integral, but their local
balances differ. At the reference, the first uses prescribed
redistribution to support constant \(k=K\); the second uses
ordinary diffusion of \cref{mawf:k-profile}.
The physical reference transport \(-Q_0\) is the same,
while the evolving models are different.

Setting the host's old switch to \(\chi=0\) does not select this
variant when that switch also restores the original-cutoff
dissipation. The present variant retains the budget cutoff
while removing both the prescribed redistribution source and
its equivalent auxiliary flux. Its scalar policy has no
spatial \(\chi\) fade. The earlier bounded-redistribution
positivity remedy would be a further model change if retained.

\subsection{A specified extension to curved terrain}
\label{mawf:curved}
Use the conservative curved momentum completion with the new
endpoint sample \(K=k|_{\Gamma_H}\), compatible coefficients,
metric-weighted reference budget and \(C_k^w=0\).
The equations are restated here so that the scalar replacement
and roughness-force prescription are explicit.

Let \(\bm n\) point from the ground into the fluid and use the
normal map
\begin{equation}
 \bm x(\bm\xi,r)=\bm X(\bm\xi)+r\bm n(\bm\xi),\qquad
 P=\Id-\bm n\otimes\bm n,\quad B_g=\nabla_\Gamma\bm n,\quad
 A_g=I_T+rB_g,\quad J=\det_TA_g.
 \label{mawf:map}
\end{equation}
The inverses of \(A_g\) act tangentially and are extended by zero
normally. The principal stretches \(1+r\gamma_i\) must be
positive, and the normal map must be globally one-to-one on
\(0\le r\le H(\bm\xi)\). The volume measure is
\(J\,\dd S_\Gamma\,\dd r\). Surface derivatives below are
taken at fixed \(r\) before restriction to the interface.

\paragraph{Reference cutoff in the volume measure.}
Set
\begin{equation}
 E_w(r)=\frac{A_{\eps,m}(K)}{\max(r,r_c)},\qquad
 p_0(r;K)=t_0(K)^2\frac{\nu_{t,A}(r;K)}{D_A(r;K)^2}.
\end{equation}
Use counterfactual zero-forcing production in
\begin{equation}
 \boxed{\int_0^H J(r)E_w(r)\,\dd r
        =\mathcal P_J:=\int_0^H J(r)p_0(r;K)\,\dd r.}
 \label{mawf:metric-budget}
\end{equation}
For \(\mathcal B(c)=A_{\eps,m}\int_0^H J/\max(r,c)\,\dd r\),
\[
 \mathcal B'(c)=-\frac{A_{\eps,m}}{c^2}\int_0^cJ\,\dd r<0,
 \qquad \mathcal B(c)\longrightarrow+\infty\quad(c\downarrow0).
\]
Hence a unique \(0<r_c\le H\) exists exactly when
\begin{equation}
 \boxed{\mathcal P_J\ge
        \frac{A_{\eps,m}}H\int_0^H J\,\dd r.}
 \label{mawf:metric-gate}
\end{equation}
At a fixed footpoint, \(J=(1+\gamma_1r)(1+\gamma_2r)\), and
\begin{equation}
 \frac{\mathcal B(c)}{A_{\eps,m}}
 =1+\ln\frac Hc
 +(\gamma_1+\gamma_2)\left(H-\frac c2\right)
 +\gamma_1\gamma_2\left(\frac{H^2}{2}-\frac{c^2}{6}\right).
 \label{mawf:metric-explicit}
\end{equation}
One monotone scalar root solve replaces the flat exponential.

\paragraph{Actual continuity, momentum, forcing and scalar transport.}
Write \(\bm U_w=\bm u_t+w\bm n\) and
\(\bm v=A_g^{-1}\bm u_t\). With Cartesian velocity components,
\begin{equation}
 \partial_r(Jw)+\operatorname{div}_\Gamma(J\bm v)=0,
 \label{mawf:continuity}
\end{equation}
\begin{equation}
 \begin{split}
 J\partial_t\bm U_w
 &+\partial_r\{J(w\bm U_w-T_w\bm n)\}\\
 &+\operatorname{Div}_\Gamma
   \{J(\bm U_w\otimes\bm v-T_wA_g^{-T})\}=J\bm f_w .
 \end{split}
 \label{mawf:curved-momentum}
\end{equation}
Here \(T_w\) is \cref{mawf:stress};
the outer stress is \(T_o=-\pi_o\Id-\R_m\).
The operator \(\operatorname{Div}_\Gamma\) takes the surface
divergence of each Cartesian row. These fluxes already contain
the geometric effects of stress and pressure.

Retain the explicit roughness-force completion
\begin{equation}
 \begin{split}
 a_r(r)&=\mathbf1_{\{r>h_r\}},\qquad
 \bm f_w=\bm f_\pi+(1-a_r)\bm F_A,\\
 \bm F_A&=\mathcal E_H^+P
 \left[\operatorname{Div}_r(\bm u_{o,t}\otimes\bm u_{o,t})
             +\nabla_r\pi_o-\bm f_\pi\right],
 \qquad \bm u_{o,t}=P\bm U_o .
 \end{split}
 \label{mawf:rough-force}
\end{equation}
Here \(\bm f_\pi\) is the remaining body acceleration after the
same hydrostatic-pressure convention used outside,
\(\nabla_rf=A_g^{-T}\nabla_\Gamma f\), and
\[
 \operatorname{Div}_r Q
 =J^{-1}\operatorname{Div}_\Gamma(JQA_g^{-T})
 \quad\text{for }Q\bm n=0.
\]
The sampler \(\mathcal E_H^+\) takes one-sided outer values after
derivatives at fixed \(r\) have been evaluated; extend the resulting
\(\bm F_A\) constantly along each normal ray.
Its work has no guaranteed drag sign and must be retained in
energy accounting.

The flat scalar reconstruction is the declared quasi-steady
parallel-flow reduction of this completion: inner storage,
normal advection, lateral viscous transport and the relevant
\(\nabla_\parallel w\) contributions are neglected, and the
remaining convection/pressure/body-force balance is replaced
by the height-constant sample \(\bm F_A\).
This also replaces the inner effective-pressure gradient by the
sampled outer mechanical-pressure gradient; the discarded
\(\nabla_\parallel(\pi_w+2k_w/3-\pi_o)\) is generally nonzero.
Under those approximations
\((D_A\bm U_\parallel')'=a_r\bm F_A\).
For the curved or unsteady patch, use the full equation
\cref{mawf:curved-momentum}, without additionally imposing that
scalar reduction. All discarded terms vanish on the homogeneous
zero-forcing reference.

The modified actual inner \(k\) equation is
\begin{equation}
 \boxed{\begin{aligned}
 J\partial_t k_w
 &+\partial_r\{J(wk_w-D_{k,w}\partial_rk_w)\}\\
 &+\operatorname{div}_\Gamma
 \{J[\bm v k_w-D_{k,w}A_g^{-1}A_g^{-T}\nabla_\Gamma k_w]\}
 =J(P_A-E_w).
 \end{aligned}}
 \label{mawf:curved-k}
\end{equation}
For incompressible flow,
\[
 P_A=2\nu_{t,A}\Smat_w:\Smat_w .
\]
The isotropic stress uses local \(k_w\), while the viscosity
uses sampled \(K\). If incompressibility is enforced only weakly,
the production contraction \(-\R_w:\nabla\bm U_w\) also contains
\(-(2k_w/3)\nabla\cdot\bm U_w\); its treatment must agree with
the discrete stress-work and pressure convention.
The base-normal tangential traction is
\begin{equation}
 \bm t_A=D_A\left[
 \partial_r\bm u_t+A_g^{-T}\nabla_\Gamma w-B_gA_g^{-1}\bm u_t
 \right].
 \label{mawf:curved-traction}
\end{equation}
Curvature, normal motion and normal pressure balance are thus
retained in the full patch.

The actual source budget remains
\[
 \int_0^HJ(P_A-E_w)\,\dd r
 =\int_0^HJ(P_A-p_0)\,\dd r,
\]
which need not vanish. A frozen reference column can be defined by
\begin{equation}
 Q_J(r)=\int_0^rJ(s)[E_w(s)-p_0(s)]\,\dd s,\qquad
 k_{\rm col}(r)=K-\int_r^H
        \frac{Q_J(s)}{J(s)D_{k,w}(s)}\,\dd s.
 \label{mawf:curved-column}
\end{equation}
Its base-normal endpoint fluxes vanish.
The one-crossing proof for \(d\ge2\nu\) extends to this column
for positive \(J,D_{k,w}\).
The stronger whole-rough-branch proof is not automatically
metric weighted. The column balances only the reference source
and normal diffusion; it is not generally an equilibrium of
the complete terrain-flow equations.

\paragraph{Actual interface geometry and shared fluxes.}
For variable \(H\), put
\begin{equation}
 \bm q_H=A_{g,H}^{-T}\nabla_\Gamma H,\qquad
 \beta_H=\sqrt{1+|\bm q_H|^2},\qquad
 \bm N_H=\frac{\bm n-\bm q_H}{\beta_H}.
 \label{mawf:interface-normal}
\end{equation}
Use common velocity and total traction,
\(\bm U_w=\bm U_o=\widehat{\bm U}\) and
\(T_w\bm N_H=T_o\bm N_H\), with one conservative momentum
flux and opposite orientations on the two sides.
The mass flux per unit ground area is
\(m_H=J_H(w-\bm u_t\cdot\bm q_H)\).
At a common \(k\) trace \(\widehat k\), match
\begin{equation}
 J_H\beta_H(\bm U k-D_k\nabla k)\cdot\bm N_H
 =m_H\widehat k-J_HD_k
 \left[\partial_rk-\bm q_H\cdot A_{g,H}^{-T}\nabla_\Gamma k\right].
 \label{mawf:curved-k-flux}
\end{equation}
Zero base-normal column flux alone does not establish zero flux
through a tilted interface. Surface differentiation of the trace
obeys
\(\nabla_\Gamma K=(\nabla_\Gamma k_w)|_H+
(\partial_rk_w)|_H\nabla_\Gamma H\).
The outer dissipation trace is \(E_w(H)=A_{\eps,m}/H\).
Full normal-traction matching replaces the flat pressure offset;
there is still no independent inner dissipation-flux prescription.

\subsection{Reference realizability, positivity, and stability limits}
\label{mawf:stability}
On a positive rough reference the smallest inner covariance
eigenvalue is
\begin{equation}
 \lambda_{\min}(\R_{w,0})
 =\frac23k_{w,0}-t_0\frac{\nu_{t,A}}{D_A}
 >K\left(\frac23-\frac1{q_m}\right)>0.
 \label{mawf:covariance}
\end{equation}
All three \(q_m\) exceed \(3/2\). Their original outer reference
covariances are also positive definite. These are reference-state
results, not global guarantees for forced or separated states.

With endpoint \(K\), coefficients, cutoff and reference production
held fixed, a wall-normal perturbation \(v=k_w-k_{w,0}\) obeys
\[
 v_t=(D_{k,w}v')',\qquad v(H)=0,\qquad D_{k,w}v'(0)=0,
\]
and
\begin{equation}
 \frac12\frac{\dd}{\dd t}\int_0^H v^2\,\dd z
 =-\int_0^H D_{k,w}(v')^2\,\dd z\le0.
 \label{mawf:scalar-stability}
\end{equation}
This conditional scalar estimate is separate from full coupled
stability.

Dynamic nonnegativity of actual \(k_w\) is not automatic:
at a flat interior zero minimum, positive endpoint \(K\), and
locally zero production and diffusion,
\[
 \partial_t k_w=-E_w(K)<0.
\]
Thus the positive rough reference does not make the evolving
sampled-\(K\) sink positivity preserving. An additional source
modification or numerical positivity device must be identified
and checked separately if introduced.

The wall changes do not alter the original outer constitutive
domains or incremental momentum response. A93 can lose covariance
realizability away from the reference, A95 retains its real-radicand
restriction, and positive covariance in E94/E95 does not ensure
positive incremental momentum diffusion. The previously established
outer constitutive restrictions therefore remain relevant.

\subsection{Illustrative references and implementation implications}
\label{mawf:verification}
Take
\begin{equation}
 \ust=0.4\ \mathrm{m\,s^{-1}},\quad
 \nu=1.5\times10^{-5}\ \mathrm{m^2\,s^{-1}},\quad
 h_r=0.1\ \mathrm m,\quad H=2\ \mathrm m,\quad z_K=z_m=H.
 \label{mawf:inputs}
\end{equation}
Using the original roughness correlation gives
\cref{mawf:numerical-table}. The same sand-grain input is retained,
so aerodynamic roughness is an output rather than a common
prescribed input.

\begin{table}[htbp]
\centering\small
\caption{Analytical rough references. \(K,k(0)\) are in
\(\mathrm{m^2\,s^{-2}}\), and lengths are in metres.}
\label{mawf:numerical-table}
\begin{tabular}{@{}llrrrrr@{}}
\toprule
Model&\(\kappa\) choice&\(K\)&\(k(0)\)&\(\Lambda\)&\(z_c\)&\(z_{0,m}\)\\
\midrule
A93&original&0.54420247&1.08630845&5.90564847&0.01480928&0.00536024\\
A95&original&0.61152007&1.22075495&5.89870963&0.01491240&0.00539286\\
E94/E95&original&0.53192103&1.11770080&5.90535573&0.01481362&0.00535834\\
\midrule
A93&common \(0.41\)&0.54420247&1.13595953&5.90420900&0.01483061&0.00536413\\
A95&common \(0.41\)&0.61152007&1.20319430&5.89920134&0.01490507&0.00539152\\
E94/E95&common \(0.41\)&0.53192103&1.12369422&5.90518519&0.01481614&0.00535880\\
\bottomrule
\end{tabular}
\end{table}

All six cases satisfy \(d\ge2\nu\) and \(\Lambda>1\), so their
positive monotone profiles have an analytical guarantee.
The interface pressure offsets are
\(-0.034324820606\), \(-0.165777590815\), and \(0\)
\(\mathrm{m^2\,s^{-2}}\), respectively.
At fixed \(\ust\) these offsets do not change with \(\kappa\),
because \(K=q_m\ust^2\) is unchanged; the inner pressure profile
does change with \(k_w(z)\).

Independent quadrature of the production and dissipation budgets,
the cumulative defect, and the closed primitive agrees with the
analytical expressions. Separately integrated source and diffusion
moments for Legendre degrees \(0,\ldots,6\) give a maximum absolute
weak residual of \(2.33\times10^{-16}\) across these cases;
closed \(k\) versus quadrature differs by at most
\(6.89\times10^{-15}\) in SI units. These are checks of the
formulas, not CFD or terrain validation.

The model-specific \(\kappa\) affects the viscosity slope,
dissipation amplitude, cutoff and inner profile. A further useful
identity is
\[
 \frac{\kappa^2[k_{w,0}(z)-K]}{u_m^2}
 =\mathcal F(az/d)-\mathcal F(aH/d).
\]
At common \(\kappa,\ust,H,z_0,\nu\), if each roughness input
is calibrated to the same \(z_0\), \cref{mawf:d-calibration}
makes \(a,d\) identical. Then the zero-forcing wall velocity,
production, sink and cumulative defect coincide across the models,
and their \(k\) profiles differ only by \(K=q_m\ust^2\).
Their forced responses can still differ through the calibrated
\(h_r\) in the forcing cutoff.

For a spectral/\(hp\) implementation, the exact continuum weak
identity \cref{mawf:weak} does not by itself make a finite polynomial
interpolant stationary. The ordinary diffusive flux, nonzero source
moments and shared interface terms must be represented and integrated
consistently. A local discrete wall solve using the same source,
quadrature and boundary contract can supply a stationary discrete
profile. Checking only the cell-mean energy budget is insufficient.

The selected variant therefore consists of endpoint sampling,
the two compatible coefficients, the budget cutoff and its plateau,
retained finite-\(\nu\) transport with \(P_A-E_w\), and conservative
velocity/stress/scalar coupling. Its exactness claim is the
flat neutral reference of these specified equations. Positivity
outside the rough theorem, global coupled stability and fully
discrete preservation remain separate requirements.
